\documentclass[dvipdfmx,a4paper,11pt,fleqn,openany]{jsbook}

\usepackage{amsmath,amsthm,amssymb,amsfonts}
\usepackage[margin=20truemm]{geometry}
\usepackage{tikz}
\usepackage{graphicx}
\usepackage{here}
\usepackage{enumitem}
\usepackage[table]{xcolor}
\usepackage{makecell}
\usepackage{multicol}
\usepackage{comment}
\usepackage{color}
\usepackage{float}
\usepackage{framed}
\usepackage{tcolorbox}
\usepackage{quotchap}

\tcbuselibrary{breakable}
\definecolor{shadecolor}{gray}{0.80}
\renewcommand{\qed}{\unskip\nobreak\quad\qedsymbol}
\raggedcolumns
% 数式上下の余白
\setlength{\abovedisplayskip}{10pt}
\setlength{\belowdisplayskip}{10pt}
\setlength{\abovedisplayshortskip}{8pt}
\setlength{\belowdisplayshortskip}{8pt}

\setlength{\columnseprule}{0.1pt}

\setcounter{tocdepth}{3}
\setcounter{secnumdepth}{4}

% 箇条書き内で式番号を付ける
\newcommand{\itemeq}[2][]{%
  \refstepcounter{equation}%
  \item
  \makebox[\linewidth]{%
    $\displaystyle #2$
    \hfill
    (\theequation)%
  }%
  \ifx\relax#1\relax\else
    \label{#1}%
  \fi
}

\title{漸化式の超体系的解説}

\begin{document}

\date{}
\author{}
\maketitle

\tableofcontents

\chapter{はじめに}
  \section{はじめに}
    \subsection{本書について}
      本書では漸化式の解法を体系的にまとめ,それぞれについて解説する.
      各解法は,いくつかの例題を通じて解説している.
      また,その解法の発想を明らかにし,行間の少ない設計を目指した.
      例題はその種の漸化式において解けなければいけないレベルのものを3題前後用意した.
      例題に対して,方針,解答,解法を付した.
      方針は解法の手がかりを,解答は結果を,解法は結果に至る過程を記した.

      本書を読む際は,机,紙,ペン,それときちんと一行一行を読む時間を用意してからページをめくってほしい.
      教科書を読むいうのは案外難しいもので,議論を丁寧に追うのは相応の苦労を伴う.
      しかしそれを経て得た知識は,演習によって強化され今後の糧となる.
      そのために,一行一行をよく読みここに書かれていることを実際に自身で試してほしい.

      本書を読むことで,与えられた漸化式がどのような解法をもって解くことができるかの判定や,必要な手立ての検討を行えるようになる.
      きちんと読みこみ,十分な演習を行えば大学受験程度の漸化式においては得意分野と言えるぐらいにはなるだろう.     

    \subsection{本書の内容と対象読者}
      本書は漸化式を一通り解けることを目標としている.
      それに際して基本的な漸化式の解法から話をする.
      最終的には難関国公立大学に出題されるような漸化式や,数学的に興味深い性質を持つ漸化式に至るまで,初歩から発展まで幅広く扱っている.
      
      対象読者は数学が好きだが,数列の分野については演習が不十分な高校生である.
      従って大部分は高校数学の初歩的な知識と最低限の計算力を前提に進める.

    \subsection{本書でのルール}
      特に断りが無い限りは,\,$i,k,l,m,n,$は1以上の整数,\,$p,q,r,s,t,u,v,x,y,z,\alpha,\beta,\gamma,\delta,\lambda$は実数として扱う.
      また,\,$\therefore\,,\because\,,\Rightarrow\,,\Leftrightarrow$は左から順に「従って」「なぜならば」「左側が成り立つならば右側が成り立つ」「同値である」を意味する.
      加えて,$\cdots$や$\vdots$は途中を省略する時に用いる.

\chapter{漸化式の解説}
  \section{漸化式の分類}
    具体的な解法について触れる前に,解法の分類を行おう.
    まず漸化式は解ける漸化式と解けない漸化式に分類される.
    そして漸化式は「決まった手順で解ける漸化式」と「特定の性質から解ける漸化式」に分類される.
    必ず解ける漸化式はさらに分類され,各分類ごとに様々な解法がある.
    ここで言う「解ける」というのは,本書内での定義であり,厳密な定義は存在しない.
    一応の基準として高校数学程度の範囲内で既知の関数に帰着するか,有限回の演算または操作により一般項を得られるような解法が存在するものを意味している.

    本書では,漸化式の基本的なパターンを説明し,次にそこからいかにしてそのパターンに帰着させるかを議論する.
    それが応用パターンの節である.

  \section{漸化式を解く前に}\label{sec:before}
    漸化式を解く前に,根底にある発想について触れておこう.
    多くの漸化式は係数などを無視して書けば$a_{n+1}=a_n$のような形をしている.
    ここに係数や数式が足されることでどんどん問題は複雑化していく.
    しかしながら不変の事実として,漸化式はある数列の2項の関係を示している.
    だから,2行の間で"きれい"な関係が成り立っていればその漸化式は解けるのである.
    この"きれい"というのは,解いていくにつれて理解されるものであるが,その多くは
    \begin{itemize}
      \item 基本4パターン（特に等比型）になっている
      \item 左辺が$n+1$に関する式,右辺が$n$に関する式になっている
      \item 左辺と右辺が同じ値
    \end{itemize}
    などである.
    最終的にこういった形を目指しながら式変形をしているということを頭の片隅に置いておくといいだろう.
    この発想は応用パターンの節で頻繁に用いられることになる.

    もちろん,一般項を予測し帰納法によって証明する方法も有効である.
    帰納法を使って証明するかどうかに関わらず,悩んだらとりあえず$a_5$ぐらいまで書き出してみる,というのは頭の中に入れておくべきことである.
    実験をするのは漸化式に限らず数学全般において重要な概念である.
  \section{基本4パターン}
    まずは多くの漸化式の最終的な終着地点になる,基本4パターンを紹介する.
    多くの可求型漸化式はこの形に持っていくことでその一般項を得る.
    その中でも特性方程式型に帰着後,等比型に帰着させることで一般項を求めることが多い.
    他の基本パターンも後の応用パターンの礎となる発想をもつため,この節は特に丁寧に取り組むことを推奨する.
    
    そして,蛇足にはなるが,一般論として「基礎」「基本」などの文言は簡単を意味しない.
    理論の根幹を成すという意味である.
    軽んじてはならない.

    \subsection{等差型 \, $a_{n+1}=a_n+d$}
      まずはもっともシンプルな漸化式であろう,等差型について説明していく.
      この漸化式は$a_1$に対して$d$がどんどんと足されていく漸化式である.
      一般項は簡単な代入計算によって容易に示せる.
      実際にやってみると
      \begin{align*}
        &a_{n+1}=a_n+d\\
        &a_{n+1}=(a_{n-1}+d)+d\\
        &a_{n+1}=\{(a_{n-2}+d)+d\}+d\\
        &\qquad \vdots\\
        &a_{n+1}=a_1+d+d+\cdots+d\\
        &a_{n+1}=a_1+nd\\
        &\therefore \quad a_n=a_1+(n-1)d
      \end{align*}
      となり一般項が示される.
      1つ前の項をどんどん代入していくのだ.
      数式の3行目から5行目周辺の式変形がすこし心配かもしれないが,$d$の個数と$a$の添字部分の和は常に$n+1$となることから考えれば分かるだろう.
      実際,$a_{n+1}$に至るまでに何度$d$が足されたかを考えているわけだから,心配であれば実際に適当な$n$で何度か実験してみるのもいい手だろう.

      実際に例題を与えよう.

      \noindent \textbf{【例題】}
      \begin{enumerate}[label=(\arabic*),parsep=3pt]
        \item $\displaystyle a_1=1,\, a_{n+1}=a_n+1$
        \item $\displaystyle a_1=\frac{1}{2},\, a_{n+1}=a_n+\frac{1}{4}$
        \item $\displaystyle a_1=-\sqrt{3},\, a_{n+1}=a_n-\sqrt{2}$
      \end{enumerate}

      \noindent\textbf{【方針】}
      \begin{enumerate}[label=(\arabic*),parsep=3pt]
        \item 等差型の一般項
              $\displaystyle a_n=a_1+(n-1)d$
              を用います.
        \item 公差は$\displaystyle d=\frac14$です.
              等差型の一般項に代入します.
        \item 公差は$\displaystyle d=-\sqrt2$です.
              負の公差の場合も同じ公式を使います.
      \end{enumerate}

      \noindent\textbf{【解答】}
      \begin{enumerate}[label=(\arabic*),parsep=3pt]
        \item $\displaystyle a_n=1+(n-1)=n$
        \item $\displaystyle a_n=\frac12+\frac{n-1}{4}
              =\frac{n+1}{4}$
        \item $\displaystyle a_n=-\sqrt3-(n-1)\sqrt2$
      \end{enumerate}

      \noindent\textbf{【解法】}\\
      解法略.
    \subsection{等比型 \, $a_{n+1}=ra_n$}
      次に説明するのは等比型である.
      一般的な解法も先ほどと同じで,\,$a_{n+1}$から順に$a_1$に向けて書き下して行くことで求めることができる.
      \begin{align*}
        &a_{n+1}=ra_n\\
        &a_{n+1}=r(ra_{n-1})\\
        &a_{n+1}=r\{r(ra_{n-2})\}\\
        &\qquad \vdots\\
        &a_{n+1}=r\cdot r \cdot \cdots \cdot ra_1\\
        &\therefore \quad a_n=r^{n-1}a_1
      \end{align*}

      \noindent \textbf{【例題】}
      \begin{enumerate}[label=(\arabic*),parsep=3pt]
        \item $\displaystyle a_1=1,\, a_{n+1}=2a_n$
        \item $\displaystyle a_1=1,\, a_{n+1}=-a_n$
        \item $\displaystyle a_1=2,\, a_{n+1}=\frac{1}{2}a_n$
      \end{enumerate}

      \noindent\textbf{【方針】}
      \begin{enumerate}[label=(\arabic*),parsep=3pt]
        \item 等比型の一般項
              $\displaystyle a_n=a_1r^{n-1}$
              を用います.
        \item 公比は$r=-1$です.
              符号が交互に変化することに注意します.
        \item 初項と公比を一般項に代入します.
      \end{enumerate}

      \noindent\textbf{【解答】}
      \begin{enumerate}[label=(\arabic*),parsep=3pt]
        \item $\displaystyle a_n=2^{n-1}$
        \item $\displaystyle a_n=(-1)^{n-1}$
        \item $\displaystyle a_n=2\left(\frac12\right)^{n-1}$
      \end{enumerate}

      \noindent\textbf{【解法】}\\
      解法略.

    \subsection{階差型 \, $a_{n+1}=a_n+f(n)$}
      階差型は先の2つに比べて一段階難易度が高い.
      しかしきちんと先程の議論を理解して,丁寧に式を追うことで理解することができるものである.
      兎にも角にも解法を見てみよう.
      \begin{align*}
        &a_{n+1}=a_n+f(n)\\
        &a_{n+1}=(a_{n-1}+f(n-1))+f(n)\\
        &a_{n+1}=\{(a_{n-2}+f(n-2))+f(n-1)\}+f(n)\\
        &\qquad \vdots\\
        &a_{n+1}=a_1+f(1)+f(2)+\cdots+f(n)\\
        &a_{n+1}=a_1+\sum_{k=1}^{n} f(k)\\
        &\therefore \quad a_n=a_1+\sum_{k=1}^{n-1} f(k)
      \end{align*}
      $f(n)$によってこの後の処理が変わるものの,漸化式の解法としてはここまでである.
      ここからはシグマの計算の能力が問われる話で,この後にやることは等差数列の総和を求めることであったり,等比数列の総和を求めることであったり,あるいは部分分数分解を用いて解くことなどが考えられる.

      \noindent \textbf{【例題】}
      \begin{enumerate}[label=(\arabic*),parsep=3pt]
        \item $\displaystyle a_1=1,\, a_{n+1}=a_n+n$
        \item $\displaystyle a_1=1,\, a_{n+1}=a_n+2^n$
        \item $\displaystyle a_1=0,\, a_{n+1}=a_n+\frac{1}{n(n+1)}$
        \item $\displaystyle a_1=1,\, a_{n+1}=a_n+(n+1)^2$
      \end{enumerate}

      \noindent\textbf{【方針】}
      \begin{enumerate}[label=(\arabic*),parsep=3pt]
        \item 最もシンプルな形の階差数列です.先程の議論を踏まえて解きましょう.
        \item もちろん,\,$f(n)$となる関数には指数関数も含まれます.何にせよやることは\,(1)と同じです.
        \item シグマは部分分数分解で処理しましょう.
        \item 勘の良い人であれば式を見た瞬間に答えがわかるかもしれませんね.これは"あの数列"の総和と同じ見た目になりませんか?.
      \end{enumerate}

      \noindent\textbf{【解答】}
      \begin{enumerate}[label=(\arabic*),parsep=3pt]
        \item $\displaystyle a_n=1+\frac{n(n-1)}{2}.$
        \item $\displaystyle a_n=2^n-1$
        \item $\displaystyle a_n=1-\frac{1}{n}$
        \item $\displaystyle a_n=\frac{1}{6}n(n+1)(2n+1)$
      \end{enumerate}

      \noindent\textbf{【解法】}
      \begin{enumerate}[label=(\arabic*),parsep=3pt]
        \item $\displaystyle a_1=1,\,a_{n+1}=a_n+n$\\
              階差型の一般項
              \[
                a_n=a_1+\sum_{k=1}^{n-1}f(k)
              \]
              を用いると,
              \begin{align*}
                a_n
                  &=1+\sum_{k=1}^{n-1}k\\
                  &=1+\frac{n(n-1)}{2}.
              \end{align*}

        \item $\displaystyle a_1=1,\,a_{n+1}=a_n+2^n$\\
              同様に,
              \begin{align*}
                a_n
                  &=1+\sum_{k=1}^{n-1}2^k\\
                  &=1+(2+2^2+\cdots+2^{n-1})\\
                  &=1+(2^n-2)\\
                  &=2^n-1.
              \end{align*}

        \item $\displaystyle a_1=0,\,a_{n+1}=a_n+\frac{1}{n(n+1)}$\\
              部分分数分解を用いると,
              \[
                \frac{1}{k(k+1)}
                =\frac{1}{k}-\frac{1}{k+1}
              \]
              であるから,
              \begin{align*}
                a_n
                  &=\sum_{k=1}^{n-1}\frac{1}{k(k+1)}\\
                  &=\sum_{k=1}^{n-1}
                    \left(\frac{1}{k}-\frac{1}{k+1}\right)\\
                  &=\left(1-\frac12\right)
                    +\left(\frac12-\frac13\right)
                    +\cdots
                    +\left(\frac{1}{n-1}-\frac1n\right)\\
                  &=1-\frac1n.
              \end{align*}

        \item $\displaystyle a_1=1,\,a_{n+1}=a_n+(n+1)^2$\\
              階差型の一般項を用いると,
              \begin{align*}
                a_n
                  &=1+\sum_{k=1}^{n-1}(k+1)^2\\
                  &=1+\sum_{j=2}^{n}j^2\\
                  &=\sum_{j=1}^{n}j^2\\
                  &=\frac{1}{6}n(n+1)(2n+1).
              \end{align*}
      \end{enumerate}
    \subsection{階比型 \, $a_{n+1}=f(n)a_n$}
      階比型漸化式は,基本4パターンの中で最も難しく,
      また,式変形が技巧的であったり,他の解法と混ざってしまいやすく,基本としての扱いが適切であるかは一考の余地がある.
      しかしながら,漸化式そのもののシンプルさの観点から基本パターンの方で扱う事とした.
      
      具体的な説明の前に1つだけ記号を新たに導入しておく.
      適当な数が並べられた数列$x_n$が存在した時,
      \[
        \prod_{i=1}^n x_i
        =x_1\times x_2\times\cdots\times x_n
      \]
      となるような記号$\prod$（ラージパイ）を定義する.
      これを総積記号といい,積を表す.
      また,一般に以下が成り立つ.
      \begin{align*}
        &\prod_{i=1}^{n}i=n!\\
        &\prod_{i=1}^{n}k=k^n\\
        &\prod_{i=1}^{n}a_ib_i=\displaystyle\prod_{i=1}^na_i\displaystyle\prod_{i=1}^nb_i\\
        &\prod_{i=1}^{n}a_i=\displaystyle\prod_{i=1+k}^{n+k}a_{i-k}\\
        &\prod_{i=1}^{n}(i+k)=\dfrac{(n+k)!}{k!}
      \end{align*}
      ここでの$a_i$や$b_i$は漸化式等によって定義されるような,法則性を持った数列ではなく
      「たくさんの値が無数に存在し,それぞれに1から順に番号が振られている集合」ぐらいのイメージで問題ない.
      では再び, \,$a_n$の一般項を求めていく.
      \begin{align*}
        &a_{n+1}=f(n)a_n\\
        &a_{n+1}=f(n)\{f(n-1)a_{n-1}\}\\
        &a_{n+1}=f(n)\left[f(n-1)\{f(n-2)a_{n-2}\}\right]\\
        &\qquad \vdots\\
        &a_{n+1}=f(1)f(2) \cdots f(n)a_1\\
        &a_{n+1}=a_1 \prod_{k=1}^n f(k)\\
        &\therefore \quad a_n=a_1 \prod_{k=1}^{n-1} f(k)
      \end{align*}
      また,以下のような解法もある.
      \begin{align*}
        &a_{n+1}=f(n)a_n\\
        &g(n+1)a_{n+1}=g(n)a_n\\
        &g(n+1)a_{n+1}=g(n)a_n=g(n-1)a_{n-1}\\
        &=\cdots=g(1)a_1\\
        &\therefore \quad a_n=\frac{g(1)a_1}{g(n)}
      \end{align*}
      この解法は,左辺と右辺を同型にしているところが肝心な部分である.
      ごちゃごちゃしたよく分からない形と同じ形を比べると,後者のほうが断然法則性が見えやすいのは説明するまでもない.
      これを背景に考えると,特に,\,$g(n)$が$n$に関する指数に$n$を持たない$m$次式のとき,例えば$f(n)=n$のとき,
      \begin{align*}
        &a_{n+1}=na_n\\
        &\frac{a_{n+1}}{n!}=\frac{a_n}{(n-1)!}\\
        &\frac{a_{n+1}}{n!}=\frac{a_n}{(n-1)!}=\frac{a_{n-1}}{(n-2)!}=\cdots=\frac{a_1}{0!}\\
        &\therefore \quad a_n=(n-1)!\cdot a_1
      \end{align*}
      が言える.
      演習問題はこの手法を使って解くものを揃えた.
      原則は左辺と右辺を同型にするということである.

      \noindent \textbf{【例題】}
      \begin{enumerate}[label=(\arabic*),parsep=3pt]
        \item $\displaystyle a_1=1,\, a_{n+1}=2na_n$
        \item $\displaystyle a_1=2,\, a_{n+1}=\frac{n+2}{n}a_n$
        \item $\displaystyle a_1=1,\, a_{n+1}=n(n+1)a_n$
      \end{enumerate}

      \noindent\textbf{【方針】}
      \begin{enumerate}[label=(\arabic*),parsep=3pt]
        \item 先程の例の類題です.
        \item 式変形をいろいろ試していくと,気付くことがあるかもしれません.大原則,左辺を$g(n+1)$に,右辺を$g(n)$に,です.
        \item ちょっと頭を捻ってみましょう.なぜ,先程までは階乗で解けて,これは無理なのでしょうか.そもそもなんで階乗で割ったのでしょうか."きれい"にしたいのでしたよね.
      \end{enumerate}

      \noindent\textbf{【解答】}
      \begin{enumerate}[label=(\arabic*),parsep=3pt]
        \item $\displaystyle a_n=(n-1)!\cdot 2^{n-1}$
        \item $\displaystyle a_n=n(n+1)$
        \item $\displaystyle a_n=n!(n-1)!$
      \end{enumerate}

      \noindent\textbf{【解法】}
      \begin{enumerate}[label=(\arabic*),parsep=3pt]
        \item $\displaystyle a_1=1,\, a_{n+1}=2na_n$\\
              両辺を$n!$で割って,
              \begin{align*}
                &a_{n+1}=2na_n\\
                &\frac{a_{n+1}}{n!}=2\frac{a_n}{(n-1)!}
              \end{align*}
              ここで
              \[
                b_n=\frac{a_n}{(n-1)!}
              \]
              と置くと,
              \begin{align*} 
                &b_{n+1}=2b_n\\
                &b_n=b_1\cdot 2^{n-1}
              \end{align*}
              $b_n$を$a_n$に戻して
              \begin{align*} 
                &\frac{a_n}{(n-1)!}=\frac{a_1}{(1-1)!}2^{n-1}\\
                &\therefore \quad a_n=a_1 \cdot (n-1)! \cdot 2^{n-1}
              \end{align*}
              
        \item $\displaystyle a_1=2,\, a_{n+1}=\frac{n+2}{n}a_n$
              \begin{align*}
                &a_{n+1}=\frac{n+2}{n}a_n\\
                &na_{n+1}=(n+2)a_n\\
              \end{align*}
              両辺を
              \[
                \displaystyle \frac{1}{n(n+1)(n+2)}
              \]
              で割って,
              \begin{align*}        
                &\frac{a_{n+1}}{(n+1)(n+2)}=\frac{a_{n}}{n(n+1)}
              \end{align*}
              $a_n$について遡っていくと
              \begin{align*}
                &\frac{a_{n}}{n(n+1)}=\frac{a_{n-1}}{(n-1)(n)}=\frac{a_{n-2}}{(n-2)(n-1)}=\cdots=\frac{a_{1}}{1(1+1)}\\
                &\therefore \quad \frac{a_{n}}{n(n+1)}=\frac{a_{1}}{1(1+1)}\\
                &a_n=n(n+1)
              \end{align*}
        \item $\displaystyle a_1=1,\, a_{n+1}=n(n+1)a_n$\\
                両辺を$n!(n+1)!$で割って,
              \begin{align*}
                &a_{n+1}=n(n+1)a_n\\
                &\frac{a_{n+1}}{n!(n+1)!}=\frac{n(n+1)a_n}{n!(n+1)!}\\
                &\frac{a_{n+1}}{n!(n+1)!}=\frac{a_n}{n!(n-1)!}\\
                &\therefore \quad \frac{a_n}{n!(n-1)!}=\frac{a_1}{1!(1-1)!}\\
                &a_n=n!(n-1)!
              \end{align*}
      \end{enumerate}

      なかなか技巧的な式変形が多く,最初は受け入れがたいかもしれない.
      しかしこういった漸化式の解法で頻出する発想は常にできるように準備しておかねばならない.

      なお,ここでは階比型漸化式$a_{n+1}=f(n)a_n$について,\,$f(n)$が指数関数でない場合について扱ったが,
      より複雑な関数であったり指数関数を用いたりする場合は対数を用いることがある.
      それら次節応用7パターンで解説する.
      
  \section{おまけ：定数数列}
    数列の種類には定数数列というものがある.
    時折,漸化式を解く中で定数数列が生まれることがある.
    先程の解説から引用しよう.
    \[
      \frac{a_{n}}{n(n+1)}=\frac{a_{n-1}}{(n-1)(n)}=\frac{a_{n-2}}{(n-2)(n-1)}=\cdots=\frac{a_{1}}{1(1+1)}
    \]
    これはまさに定数数列の良い例である.
    少し式が複雑なので,
    \[
      b_n=\frac{a_{n}}{n(n+1)}
    \]
    とおくと式は
    \[
      b_n=b_{n-1}=b_{n-2}=\cdots=b_1
    \]
    となる.
    ずっと同じ値を取るのだ.
    当たり前過ぎて何を言っているんだ,と思うかもしれないが急に出てきてパニックにならないようにここで念のため触れておく.

\section{応用7パターン}
  ここからは,いかにして基本4パターンに帰着させるか,という話をしていく.
  つまりは \ref{sec:before} でいう"きれい"な形を目指すということである.
  問題を解くときには,いかにして既知の形にするか,いかにして単純な形にするかを考える.
  漸化式で言えば,基本4パターンがそれに該当する.
  この形であれば解けるというのを1つづつ増やしていき,それらを積み重ねる.
  漸化式,ひいては数学において重要なことである.
  
  \subsection{特性方程式型 \, $a_{n+1}=pa_n+q$}
    応用パターンの1つ目,特性方程式型について説明する.
    まずは一般的な形を見せよう.
    \[
      a_{n+1}=pa_n+q
    \]
    さて,この漸化式はどのように解くことができるだろう.
    漸化式をうまく変形して知っている形に変形できれば解くことができる.
    先に扱った基本4パターンの知識を使ってこの漸化式を分析してみよう.
    $a_n$に対して$p$が掛けられ,さらに$q$が足されている.
    つまり,この漸化式は,等比型と階差型の組み合わせからなっていることがわかる.
    そこから判断して,ここは少々突飛な発想ではあるが,この漸化式を等比型に帰着させることを試みよう.
    
    この漸化式は,等比型漸化式
    \[
      a_{n+1}=pa_n
    \]
    に対して$q$が足されたものだ.
    つまりは,\,$q$がうまく消えてくれるような計算式を考えたいのだ.
    その分だけうまく戻せばよい.
    結論から言えば,
    \[
      \begin{array}{rrl}
         & a_{n+1}={} & pa_n+q\\
        -\,) & \alpha={} & p\alpha+q\\ \hline
         & a_{n+1}-\alpha={} & p(a_n-\alpha)
      \end{array}
    \]
    より,
    \[
      a_{n+1}-\alpha=p(a_n-\alpha)
    \]
    を満たすような$\alpha$を考えればよい.
    この式を特性方程式と言う.
    形式的には,
    \[
      a_{n+1}-\alpha=p(a_n-\alpha)
    \]
    または
    \[
      \alpha=p\alpha+q
    \]
    を解き,\,$\alpha$を求めればよい.
    続けると,
    \begin{align*}
      &a_{n+1}-\alpha=p(a_n-\alpha)\\
      &a_{n+1}-\alpha=pa_n-p\alpha\\
      &a_{n+1}=pa_n-p\alpha+\alpha\\
      &a_{n+1}=pa_n+(1-p)\alpha
    \end{align*}
    $(1-p)\alpha$について,
    \[
      a_{n+1}=pa_n+q
    \]
    と係数比較を行うと,
    \begin{align*}
      &q=(1-p)\alpha\\
      &\therefore \quad \alpha=\frac{q}{1-p}
    \end{align*}
    $\alpha$が求まったので,次のステップに移る.
    式の見た目を整えるために,\,$b_n=a_n-\alpha$と置く.
    すると,
    \begin{align*}
      &a_{n+1}-\alpha=p(a_n-\alpha) \iff b_{n+1}=pb_n
    \end{align*}
    この等比型漸化式$b_{n+1}=pb_n$を解くと,
    \begin{align*}
      &b_{n+1}=pb_n\\
      &\therefore \quad b_n=b_1p^{n-1}\\
      &\therefore \quad a_n-\alpha=(a_1-\alpha)p^{n-1}\\
      &\therefore \quad a_n=(a_1-\alpha)p^{n-1}+\alpha\\
      &\therefore \quad a_n=\left(a_1-\frac{q}{1-p}\right)p^{n-1}+\frac{q}{1-p}
    \end{align*}
    となり,\,$a_n$の一般項が示された.
    やってほしくないことは,この議論の末に得られた結果の形だけを覚えて問題に挑むことである.
    この議論全体を問題を解くたびに都度再現してほしい.

    もう少し数学の世界に深入りしていこう.
    \,$\alpha$を求めるというのはどういった意味合いを持つだろうか.
    一度,この漸化式を$y=px+q$とみなしてみよう.
    すると,この漸化式は$y=px$の$y$軸方向の平行移動と捉えることができる.
    これはすなわち,
    \[
      x_{n+1}=f(x_n)
    \]
    というような関数を考えることに相当する.
    この時に,
    \[
      x^*=f(x^*)
    \]
    を満たすような$x^*$が存在する場合を考える.
    このような点を探す行為は座標平面上において,\,$y=x$と$y=px+q$の交点を探すことに相当する.
    この交点を原点と思ってグラフをみると,この座標平面は$y=px$と$y=x$のグラフが原点から始まって進んでいくようなものに見える.
    このような点を不動点という.
    そしてそれを$\alpha$として我々は先程計算を行ったのだ.
    この不動点を知るモチベーションは,複雑な値の変化をより簡素な値の変化に置き換えることで,より深く現象を考察できるようになることである.
    つまり,不動点は特性方程式型の漸化式の複雑な値の変化を,よりシンプルな形にする役割を持つ.
    不動点を知ることで,そこへ近づく,あるいは遠ざかる値の変化を克明に捉える事ができ,それが$y=px+q$と$y=x$のグラフとの交点に該当する.
    この点においては両辺が等しいので,この点においては値の変化は起こらない.
    この話から,漸化式の作る値はの概形は,\,$p$が作るということが分かる.

    また,特性方程式型漸化式には別の解法が存在する.
    以下に示す.
    \begin{align}
      &a_{n+1}=pa_n+q \tag{1}\\
      &a_{n+2}=pa_{n+1}+q \tag{2}\\
      &\text{$(2)-(1)$をして,} \notag \\
      &a_{n+2}-a_{n+1}=p(a_{n+1}-a_n) \notag \\
      &\text{ここで,\,$b_n=a_{n+1}-a_n$と置くと,} \notag \\
      &b_{n+1}=pb_n \notag
    \end{align}
    このようにすることで,漸化式は等比数列に帰着する.
    この後の処理は読者に委ねる.

    今後,特性方程式を用いた解法を多数扱う.
    特性方程式は漸化式を解く中で重要な概念なのである.
    その背景には行列や線形代数が深く絡んでいる.

  \noindent \textbf{【例題】}
  \begin{enumerate}[label=(\arabic*),parsep=3pt]
    \item $\displaystyle a_1=1,\, a_{n+1}=2a_n+1$
    \item $\displaystyle a_1=1,\, a_{n+1}=\frac{1}{2}a_n+2$
    \item $\displaystyle a_1=1,\, a_{n+1}=a_n+2$
  \end{enumerate}

  \noindent\textbf{【方針】}
  \begin{enumerate}[label=(\arabic*),parsep=3pt]
    \item まずは,\,$\alpha$を求めて等比型に帰着させる方法で解きましょう.
    \item 同じく,\,$\alpha$を求めて等比型に帰着させる方法で解きましょう.
    \item この漸化式は等差型ですが,試しに特性方程式型に帰着させることを考えてみましょう.$\alpha$は存在しますか?
  \end{enumerate}

  \noindent\textbf{【解答】}
  \begin{enumerate}[label=(\arabic*),parsep=3pt]
    \item $\displaystyle a_n=2^n-1$
    \item $\displaystyle a_n=4-3\cdot\left(\frac{1}{2}\right)^{n-1}$
    \item $\displaystyle a_n=1+2(n-1)$
  \end{enumerate}

  \noindent\textbf{【解法】}
  \begin{enumerate}[label=(\arabic*),parsep=3pt]
    \item $\displaystyle a_1=1,\,a_{n+1}=2a_n+1$\\
          特性方程式を考えると,
          \[
            \alpha=2\alpha+1
            \quad\Longleftrightarrow\quad
            \alpha=-1
          \]
          である.

          $b_n=a_n-\alpha=a_n+1$とおくと,
          \begin{align*}
            b_{n+1}
              &=a_{n+1}+1\\
              &=2a_n+2\\
              &=2(a_n+1)\\
              &=2b_n.
          \end{align*}
          また,\,$b_1=a_1+1=2$であるから,
          \[
            b_n=2^n.
          \]
          よって,
          \[
            a_n=b_n-1=2^n-1.
          \]

    \item $\displaystyle a_1=1,\,a_{n+1}=\frac12a_n+2$\\
            特性方程式を考えると,
            \[
              \alpha=\frac12\alpha+2
              \quad\Longleftrightarrow\quad
              \alpha=4
            \]
            である.

            $b_n=a_n-\alpha=a_n-4$とおくと,
            \begin{align*}
              b_{n+1}
                &=a_{n+1}-4\\
                &=\frac12a_n+2-4\\
                &=\frac12(a_n-4)\\
                &=\frac12b_n.
            \end{align*}
            また,\,$b_1=a_1-4=-3$であるから,
            \[
              b_n=-3\left(\frac12\right)^{n-1}.
            \]
            よって,
            \[
              a_n=b_n+4
              =4-3\left(\frac12\right)^{n-1}.
            \]
     \item $\displaystyle a_1=1,\, a_{n+1}=a_n+2$\\
           $\alpha$は存在しないので,特性方程式型に帰着させることはできない.
           これが意味することは,不動点を持たないということである.
  \end{enumerate}

  \subsection{階差等比複合型 \, $\displaystyle a_{n+1}=pa_n+f(n)$}
    \subsubsection{一般型 \, $\displaystyle a_{n+1}=pa_n+f(n)$}
      では次に,階差等比複合型について説明する.
      この漸化式は,等比型と階差型の組み合わせからなっている.
      先程の特性方程式型と同じように,等比型に帰着させることを試みよう.
      \[
        a_{n+1}+g(n+1)=p\{a_n+g(n)\}
      \]
      となってくれれば,漸化式は等比型に帰着にする.
      このような$g(n)$が存在するかどうかは分からないが,一度存在を仮定して計算を進めよう.
      ここで,\,$g(n)$の次数は$f(n)$の次数に依存する.
      したがって,\,$f(n)$が$m$次式の時,$g(n)$は$m$次式である.
      この時$g(n)$は
      \[
        g(n) = x_0 + x_1 n^1 + x_2 n^2  + \cdots +x_{m} n^m =\sum_{i=0}^{m} x_i n^{i}
      \]
      のように書くことができる.
      ここで$x_i$は適当な実数が並べられた数の集合である.
      $g(n)$が求まったら,\,$b_n=a_n+g(n)$と置くと,
      \begin{align*}
        b_{n+1}=pb_n
      \end{align*}
      のように等比数列になる.
      あとは,先程の特性方程式型と同じようにして一般項を求めることができる.
    \subsubsection{指数型  \, $\displaystyle a_{n+1}=pa_n+q^n$}
      前節では$f(n)$が任意の関数である場合を扱った.
      ここでは$f(n)$が特に指数関数である場合を扱う.
      指数関数である場合にはもう一つ別の解法が存在する.
      \begin{align*}
        a_{n+1}=pa_n+q^n
      \end{align*}
      両辺を$q^n$で割って,
      \begin{align*}      
        &\frac{a_{n+1}}{q^n}=\frac{p}{q}\cdot\frac{a_n}{q^{n-1}}+1\\
        &\text{ここで,}\,b_n=\frac{a_n}{q^{n-1}}\text{と置くと,}\\
        &b_{n+1}=\frac{p}{q}b_n+1
      \end{align*}
      あとは特性方程式型の漸化式として解くことになる.
  \subsection{対数型 \, $a_{n+1}=pa_n^q$}
    今までは$a_n$に対する操作は定数倍したり関数が掛けられたりするだけだった.
    しかし,\,$a_n$に指数がついたらどうだろうか.
    この場合はかなり難しい発想の変形が必要である.
    それは対数を取るという操作である.
    実際にやってみよう.
    \[
      a_{n+1}=pa_n^q
    \]
    について,両辺の対数を取ると
    \begin{align*}
      \log_p{a_{n+1}}&=\log_p{pa_n^q}\\
      &=\log_p{p}+\log_p{a_n^q}\\
      &=q\log_p{a_n}+1
    \end{align*}
    $b_n=\log_p{a_n}$とおくと,
    \begin{align*}
      b_{n+1}=qb_n+1
    \end{align*}
    あとはただの特性方程式型の問題である.
    底をいくらにするかは自由だが$p$に関連する数だと扱いやすい.
    たとえば$p=k^m$の時は$k$でとっておくことが多い.
    作問者側は解答がシンプルな形になるようにするため,より小さい底を取ることで答えが求めやすくなる.
    \noindent\textbf{【例題】}
    \begin{enumerate}[label=(\arabic*),parsep=3pt]
      \item $\displaystyle a_1=2,\quad a_{n+1}=2a_n^2$
      \item $\displaystyle a_1=3,\quad a_{n+1}=9a_n^2$
      \item $\displaystyle a_1=1,\quad a_{n+1}=(n+1)a_n$
    \end{enumerate}

    \noindent\textbf{【方針】}
    \begin{enumerate}[label=(\arabic*),parsep=3pt]
      \item 両辺の底を$2$とする対数を取り,$b_n=\log_2a_n$とおいて一次の漸化式に変形します.その結果,漸化式は特性方程式型になります.
      \item $9=3^2$であることに注目し,底を$3$とする対数を取りましょう.あとは(1)と同様です.
      \item 階比型に見えますが,対数型としても解けます.
    \end{enumerate}

    \noindent\textbf{【解答】}
    \begin{enumerate}[label=(\arabic*),parsep=3pt]
      \item $\displaystyle a_n=2^{3\cdot2^{n-1}-2}$
      \item $a_n=3^{3\cdot2^{n-1}-2}$
      \item $a_n=n!$
    \end{enumerate}

    \noindent\textbf{【解法】}
    \begin{enumerate}[label=(\arabic*),parsep=3pt]
      \item $\displaystyle a_1=2,\quad a_{n+1}=2a_n^2$\\
            両辺の底を$2$とする対数を取ると,
            \begin{align*}
              \log_2a_{n+1}
                &=\log_2{2a_n^2}\\
                &=2\log_2{2a_n}\\
                &=2(\log_2{2}+\log_2{a_n})\\
                &=2\log_2{a_n}+2
            \end{align*}
            ここで
            \[
              b_n=\log_2a_n
            \]
            とおくと,
            \begin{align*}
              &b_{n+1}=2b_n+2\\
              &b_{n+1}+2=2(b_n+2)
            \end{align*}
            ここで
            \[
              c_n=b_n+2
            \]
            とおくと,
            \[
              c_{n+1}=2c_n
            \]
            $c_1$は,
            \[
              c_1=b_1+2=\log_2{a_1}+2=\log_2{2}+2=1+2=3
            \]
            従って,
            \begin{align*}
              &c_n=3\cdot2^{n-1}\\
              &b_n+2=3\cdot2^{n-1}\\
              &b_n=3\cdot2^{n-1}-2\\
              &a_n=2^{3\cdot2^{n-1}-2}
            \end{align*}
      \item $\displaystyle a_1=3,\quad a_{n+1}=9a_n^2$
            $9=3^2$\\
            底を$3$とする対数を取ると,
            \begin{align*}
              \log_3a_{n+1}
                =\log_3(9a_n^2)=2+2\log_3a_n
            \end{align*}
            ここで
            \[
              b_n=\log_3a_n
            \]
            とおくと,
            \begin{align*}
              &b_{n+1}=2b_n+2\\
              &b_{n+1}+2=2(b_n+2)
            \end{align*}
            ここで
            \[
              c_n=b_n+2
            \]
            とおくと,
            \[
              c_{n+1}=2c_n
            \]
            $c_1=3$より,
            \begin{align*}
              &c_n=3\cdot2^{n-1}\\
              &b_n+2=3\cdot2^{n-1}\\
              &a_n=3^{3\cdot2^{n-1}-2}
            \end{align*}
      \item $\displaystyle a_1=1,\quad a_{n+1}=(n+1)a_n$\\
            両辺の自然対数を取って
            \[
              \log a_{n+1}=\log{n+1}+\log a_n
            \]
            したがって
            \begin{align*}
              \log a_n&=\log a_1 + \log{2}+\log{3}+\cdots+ \log{n}\\
              &=\log 1 + \log{2}+\log{3}+\cdots +\log{n}\\
              &=\log n!\\
              \therefore \quad a_n=n!
            \end{align*}
    \end{enumerate}

  \subsection{分数型}
    \subsubsection{逆数型  \, $\displaystyle a_{n+1}=\frac{pa_n}{qa_n+r}$}
      では逆数型の漸化式について扱っていこう.
      この漸化式の解法はかなり技巧的な式変形が要求される.
      実際に見せよう.
      \begin{align*}
        &a_{n+1}=\frac{pa_n}{qa_n+r}\\
        &\text{両辺の逆数を取って}\\
        &\frac{1}{a_{n+1}}=\frac{qa_n+r}{pa_n}\\
        &\frac{1}{a_{n+1}}=\frac{q}{p}+\frac{r}{p}\cdot\frac{1}{a_n}\\
        &\text{$\displaystyle \frac{1}{a_n}=b_n$と置くと}\\
        &b_{n+1}=\frac{r}{p}b_n+\frac{q}{p}
      \end{align*}
      このようにすることで,特性方程式型の漸化式になる.

      \noindent \textbf{【例題】}
      \begin{enumerate}[label=(\arabic*),parsep=3pt]
        \item $\displaystyle a_1=1,\, a_{n+1}=\frac{a_n}{a_n+1}$
        \item $\displaystyle a_1=1,\, a_{n+1}=\frac{-a_n}{2a_n+1}$
      \end{enumerate}

      \noindent\textbf{【方針】}
      \begin{enumerate}[label=(\arabic*),parsep=3pt]
        \item とにかく基本に忠実に解きましょう.
        \item 基本に忠実に.
      \end{enumerate}

      \noindent\textbf{【解答】}
      \begin{enumerate}[label=(\arabic*),parsep=3pt]
        \item $\displaystyle a_n=\frac{1}{n}$
        \item $\displaystyle a_n=\frac{1}{2(-1)^{n-1}-1}$
      \end{enumerate}

      \noindent\textbf{【解法】}
      \begin{enumerate}[label=(\arabic*),parsep=3pt]
        \item $\displaystyle a_1=1,\, a_{n+1}=\frac{a_n}{a_n+1}$\\
              両辺の逆数を取って,
              \begin{align*}
                &\frac{1}{a_{n+1}}=\frac{a_n+1}{a_n}=1+\frac{1}{a_n}\\
                &\text{$b_n=\dfrac{1}{a_n}$と置くと,}\\
                &b_{n+1}=b_n+1,\quad b_1=\frac{1}{a_1}=1\\
                &\therefore \quad b_n=1+(n-1)\cdot1=n\\
                &\therefore \quad a_n=\frac{1}{b_n}=\frac{1}{n}
              \end{align*}
        \item $\displaystyle a_1=1,\, a_{n+1}=\frac{-a_n}{2a_n+1}$\\
              両辺の逆数を取って,
              \begin{align*}
                &\frac{1}{a_{n+1}}=\frac{2a_n+1}{-a_n}=-2-\frac{1}{a_n}\\
                &\text{$b_n=\dfrac{1}{a_n}$と置くと,}\\
                &b_{n+1}=-b_n-2,\quad b_1=1
              \end{align*}
              特性方程式型として$\alpha$を求めると,
              \[
                \alpha=-\alpha-2 \ \Leftrightarrow\ \alpha=-1
              \]
              なので,\,$c_n=b_n+1$と置くと,
              \begin{align*}
                &c_{n+1}=-c_n,\quad c_1=b_1+1=2\\
                &\therefore \quad c_n=2\cdot(-1)^{n-1}\\
                &\therefore \quad b_n=c_n-1=2(-1)^{n-1}-1\\
                &\therefore \quad a_n=\frac{1}{b_n}=\frac{1}{2(-1)^{n-1}-1}
              \end{align*}
      \end{enumerate}
    \subsubsection{1次分数型 \, $\displaystyle a_{n+1}=\frac{pa_n+q}{ra_n+s}$}
      先程の逆数型の一般化と言える漸化式である.
      ここで一度断りを入れておこう.この小節では,分母が0にならないものとして議論を進める.
      0になる場合は,そもそも数列が定義できない,あるいは,別の漸化式に帰着するからである.

      数学における基本的な指針として,既知の形に帰着させるというものがある.
      本書で言う,"きれい"な形にするということである.
      今回の場合,もっとも一次分数型に近い漸化式は逆数型であることから,逆数型に帰着させることで解けると予想しよう.
      \[
        a_{n+1}=\frac{pa_n+q}{ra_n+s} \Leftrightarrow b_{n+1}=\frac{tb_n}{u b_n+v} 
      \]
      となるような$b_n$と$t,u,v$の組があればよい.
      ここでも特性方程式型のアイデアは活きる.
      この時$b_n$は$a_n$について$f(n)$の形を含まないから\,$b_n=a_n-\alpha$と置くことができる. 
      また,\,$b_n=a_n-\alpha$から$a_n=b_n+\alpha$が言える.
      ひとまずこの$\alpha$を求めていこう.
      そのために
      \[
        a_{n+1}=\frac{pa_n+q}{ra_n+s} \Leftrightarrow a_{n+1}=\frac{p(b_n+\alpha)+q}{r(b_n+\alpha)+s}
      \]
      の形に変形できると仮定して,計算を進める.
      \begin{align*}
        &a_{n+1}=\frac{p(b_n+\alpha)+q}{r(b_n+\alpha)+s}\\
        &b_{n+1}+\alpha=\frac{p(b_n+\alpha)+q}{r(b_n+\alpha)+s}\\
        &b_{n+1}=\frac{p(b_n+\alpha)+q}{r(b_n+\alpha)+s}-\alpha\\
        &b_{n+1}=\frac{p(b_n+\alpha)+q-\alpha\{r(b_n+\alpha)+s\}}{r(b_n+\alpha)+s}\\
        &b_{n+1}=\dfrac{(p-r\alpha)b_{n}+p\alpha+q-r\alpha^{2}-s\alpha}{rb_{n}+r\alpha+s}\\
      \end{align*}
      ここで,\,$p\alpha+q-r\alpha^{2}-s\alpha=0$であれば,目標であった
      \[
        a_{n+1}=\frac{pa_n+q}{ra_n+s} \Leftrightarrow b_{n+1}=\frac{tb_n}{u b_n+v} 
      \]
      の形の達成である.また,
      \[
        p\alpha+q-r\alpha^{2}-s\alpha=0 \Leftrightarrow \alpha=\frac{p\alpha+q}{r\alpha+s}
      \]
      であることから,形式的に$a_n$を$\alpha$に置き換えた,\,$\alpha$についての二次方程式を解くことで求めることができる.
      あとは逆数型として計算してもよいのだが,計算が煩雑であるから,より良い方法について検討しよう.

      ここでこの話は再び不動点に関しての議論に至る.
  
      ここまででは,\,$b_n=a_n-\alpha$と置くことで,
      \[
        b_{n+1}=\frac{tb_n}{ub_n+v}
      \]
      の形に帰着させることを考えた.
      
      ここで,もう少しこの$\alpha$について考えてみよう.
      先ほどの計算では
      \[
        p\alpha+q-r\alpha^2-s\alpha=0
      \]
      となるように$\alpha$を選んだ.
      これは整理すると
      \[
        r\alpha^2+(s-p)\alpha-q=0
      \]
      である.
      
      つまり,\,$\alpha$は適当に選んだ数ではなく,
      \[
        rx^2+(s-p)x-q=0
      \]
      という二次方程式の解なのである.
      
      ここで,この二次方程式が異なる2つの解$\alpha,\beta$を持つ場合を考えよう.
      ここでの$\alpha,\beta$は先程の$\alpha$が取りうる実数の2解を別々に書いたものと言う認識で問題ない.
      解と係数の関係より,
      \[
        \alpha+\beta=\frac{p-s}{r},
        \qquad
        \alpha\beta=-\frac qr
      \]
      である.
      したがって,$\alpha,\beta$を個別に解の公式で求めなくても,
      \[
        x^2-(\alpha+\beta)x+\alpha\beta
        =
        x^2-\frac{p-s}{r}x-\frac qr
      \]
      と書くことができる.

      この二次式が重要なのは,$\alpha,\beta$を用いて
      \[
        \frac{a_{n+1}-\alpha}{a_{n+1}-\beta}
      \]
      という形を作ることができるからである.

      実際,
      \[
        a_{n+1}=\frac{pa_n+q}{ra_n+s}
      \]
      を代入すると,
      \begin{align*}
        \frac{a_{n+1}-\alpha}{a_{n+1}-\beta}
        &=
        \frac{(p-r\alpha)a_n+q-s\alpha}{(p-r\beta)a_n+q-s\beta}
      \end{align*}
      
      ここで$\alpha$は
      \[
        r\alpha^2+(s-p)\alpha-q=0
      \]
      を満たすので,
      \[
        q-s\alpha
        =\alpha(r\alpha+s-p)
      \]
      よって,
      \[
        (p-r\alpha)a_n+q-s\alpha
        =(p-r\alpha)(a_n-\alpha)
      \]
      $\beta$についても同様に,
      \[
        (p-r\beta)a_n+q-s\beta
        =(p-r\beta)(a_n-\beta)
      \]
      したがって,
      \begin{align*}
        \frac{a_{n+1}-\alpha}{a_{n+1}-\beta}
        &=
        \frac{p-r\alpha}{p-r\beta}
        \frac{a_n-\alpha}{a_n-\beta}
      \end{align*}
      
      ここで
      \[
        b_n=\frac{a_n-\alpha}{a_n-\beta}
      \]
      と置けば,
      \[
        b_{n+1}
        =
        \frac{p-r\alpha}{p-r\beta}b_n
      \]
      となり,等比数列が得られる.
      $\alpha$と$\beta$の大小関係は,解の大小関係を要求されていないことから自由にとって良い.
      つまり,\,1次分数型漸化式では,
      \[
        a_{n+1}=\frac{pa_n+q}{ra_n+s}
      \]
      という複雑な形をそのまま扱うのではなく,
      \[
        \frac{a_n-\alpha}{a_n-\beta}
      \]
      という「$\alpha,\beta$からの位置の比」に置き換えることで,
      等比数列に変えることができる.

      そして$\alpha,\beta$は,
      \[
        rx^2+(s-p)x-q=0
      \]
      の2解として自然に現れる.
      
      したがって$\alpha$は,もとの漸化式の$a_n$を形式的に$\alpha$に置き換えたときにも値が変化しない数,すなわち不動点である.
      $\beta$についても同じことがいえる.
      
      議論が長くなったので,ここで解法をまとめておこう.
      \[
        a_{n+1}=\frac{pa_n+q}{ra_n+s} 
      \]
      の形で表される漸化式に対して,
      \[
        x= \frac{px+q}{rx+s} \Leftrightarrow rx^2+(s-p)x-q=0
      \]
      を解き,その2解を$\alpha,\beta$とする.
      $r\neq0$であれば,
      \[
        x^2+\frac{s-p}{r}x-\frac qr=0
      \]
      と書けるので,解と係数の関係から
      \[
        \alpha+\beta=\frac{p-s}{r},
        \qquad
        \alpha\beta=-\frac qr
      \]
      である.
      ここで,
      \[
        b_n=\frac{a_n-\alpha}{a_n-\beta}
      \]
      と置く.
      すると,
      \begin{align*}
        b_{n+1}&=\frac{a_{n+1}-\alpha}{a_{n+1}-\beta}\\
        &=\frac{\displaystyle\frac{pa_n+q}{ra_n+s}-\alpha}{\displaystyle\frac{pa_n+q}{ra_n+s}-\beta}\\
        &=\frac{pa_n+q-\alpha(ra_n+s)}{pa_n+q-\beta(ra_n+s)}\\
        &=\frac{(p-r\alpha)a_n+q-s\alpha}{(p-r\beta)a_n+q-s\beta}
      \end{align*}
      ここで,$\alpha$は
      \[
        r\alpha^2+(s-p)\alpha-q=0
      \]
      を満たすから,
      \[
        q-s\alpha=-\alpha(p-r\alpha)
      \]
      である.
      したがって,
      \[
        (p-r\alpha)a_n+q-s\alpha
        =(p-r\alpha)(a_n-\alpha).
      \]
      同様に,$\beta$について
      \[
        (p-r\beta)a_n+q-s\beta
        =(p-r\beta)(a_n-\beta)
      \]
      である.
      よって,
      \begin{align*}
        b_{n+1}&=\frac{(p-r\alpha)(a_n-\alpha)}{(p-r\beta)(a_n-\beta)}\\
        &=\frac{p-r\alpha}{p-r\beta}\frac{a_n-\alpha}{a_n-\beta}\\
        &=\frac{p-r\alpha}{p-r\beta}b_n
      \end{align*}
      したがって,
      \[
        b_{n+1}=\frac{p-r\alpha}{p-r\beta}b_n
      \]
      となり,$b_n$は等比数列である.
      初項
      \[
        b_1=\frac{a_1-\alpha}{a_1-\beta}
      \]
      を用いれば,
      \[
        b_n=\frac{a_1-\alpha}{a_1-\beta}\left(\frac{p-r\alpha}{p-r\beta}\right)^{n-1} \tag{＊}
      \]
      となる.
      最後に,
      \[
        b_n=\frac{a_n-\alpha}{a_n-\beta}
      \]
      から$a_n$を求める.
      \begin{align*}
        b_n(a_n-\beta)&=a_n-\alpha\\
        (b_n-1)a_n&=\beta b_n-\alpha\\
        a_n=\frac{\beta b_n-\alpha}{b_n-1}
      \end{align*}
      (＊)を代入して$a_n$を求めることができる.

      \noindent \textbf{【例題】}
      \begin{enumerate}[label=(\arabic*),parsep=3pt]
        \item $\displaystyle a_1=3,\, a_{n+1}=\frac{2a_n+1}{a_n+2}$
        \item $\displaystyle a_1=0,\, a_{n+1}=\frac{a_n+4}{a_n+1}$
        \item $\displaystyle a_1=5,\, a_{n+1}=\dfrac{4a_{n}-9}{a_{n}-2}$
        \item $\displaystyle a_1=1,\, a_{n+1}=1+\frac{1}{a_n}$
      \end{enumerate}

      \noindent\textbf{【方針】}
      \begin{enumerate}[label=(\arabic*),parsep=3pt]
        \item 最も基本的な例です.先程の議論に習って解きましょう.
        \item 同様に不動点$\alpha,\beta$を求めましょう.今回は片方が負の値になります.
        \item 今回は不動点の方程式が重解になります.それでもやることは同じです.
        \item このような形も一次分数型です.解と係数の関係を活用して解きましょう.
      \end{enumerate}

      \noindent\textbf{【解答】}
      \begin{enumerate}[label=(\arabic*),parsep=3pt]
        \item $\displaystyle a_n=\frac{2\cdot3^{n-1}+1}{2\cdot3^{n-1}-1}$
        \item $\displaystyle a_n=\frac{2\left(3^{n-1}+(-1)^n\right)}{3^{n-1}-(-1)^n}$
        \item $a_{n}=\dfrac{6n-1}{2n-1}$
        \item $\displaystyle a_n=\cfrac{\left(\cfrac{1-\sqrt{5}}{2}\right)^{n+1}-\left(\cfrac{1+\sqrt{5}}{2}\right)^{n+1}}{\left(\cfrac{1-\sqrt{5}}{2}\right)^n-\left(\cfrac{1+\sqrt{5}}{2}\right)^n}$
      \end{enumerate}

      \noindent\textbf{【解法】}
      \begin{enumerate}[label=(\arabic*),parsep=3pt]
        \item $\displaystyle a_1=3,\, a_{n+1}=\frac{2a_n+1}{a_n+2}$\\
              不動点の方程式は
              \[
                x=\frac{2x+1}{x+2}
              \]
              であるから,その解をそれぞれ$\alpha,\beta$とすると\,$\alpha=1,\ \beta=-1$と取れる.
              このとき,
              \[
                b_n=\frac{a_n-1}{a_n+1}
              \]
              と置くと,
              \begin{align*}
                b_{n+1}&=\frac{a_{n+1}-1}{a_{n+1}+1}\\
                &=\cfrac{\cfrac{2a_n+1}{a_n+2}-1}{\cfrac{2a_n+1}{a_n+2}+1}\\
                &=\frac{2a_n+1-(a_n+2)}{2a_n+1+(a_n+2)}\\
                &=\frac{a_n-1}{3a_n+3}\\
                &=\frac{1}{3}\cdot \frac{a_n-1}{a_n+1}
              \end{align*}
              $\displaystyle b_1=\frac{a_1-1}{a_1+1}=\frac{2}{4}=\frac{1}{2}$から,
              \begin{align*}
                b_{n}=\frac{1}{2}\cdot\left(\frac{1}{3}\right)^{n-1}\\
                \therefore \quad b_n=\frac{1}{2\cdot3^{n-1}}
              \end{align*}
              最後に
              \[
                b_n=\frac{a_n-1}{a_n+1}
              \]
              を$a_n$について解くと,
              \begin{align*}
                &b_n(a_n+1)=a_n-1\\
                &(b_n-1)a_n=-1-b_n\\
                &a_n=\frac{1+b_n}{1-b_n}
              \end{align*}
              \[
                b_n=\frac{1}{2\cdot3^{n-1}}
              \]
              を代入して整理すれば,
              \[
                a_n=\frac{2\cdot3^{n-1}+1}{2\cdot3^{n-1}-1}
              \]
              を得る.
        \item $\displaystyle a_1=0,\, a_{n+1}=\frac{a_n+4}{a_n+1}$\\
              不動点の方程式は
              \[
                x=\frac{x+4}{x+1}
              \]
              であるから,その解をそれぞれ$\alpha,\beta$とすると\,$\alpha=2,\ \beta=-2$と取れる.
              このとき,
              \[
                \frac{p-r\alpha}{p-r\beta}=\frac{1-2}{1+2}=-\frac13
              \]
              なので,
              \[
                b_n=\frac{a_n-2}{a_n+2}
              \]
              と置くと,
              \begin{align*}
                &b_{n+1}=-\frac13 b_n,\quad b_1=\frac{a_1-2}{a_1+2}=\frac{-2}{2}=-1\\
                &\therefore \quad b_n=(-1)\left(-\frac13\right)^{n-1}=\frac{(-1)^n}{3^{n-1}}
              \end{align*}
              \[
                b_n=\frac{a_n-2}{a_n+2}
              \]
              を$a_n$について解くと,
              \[
                a_n=\frac{2(1+b_n)}{1-b_n}
              \]
              となるので,
              \[
                b_n=\frac{(-1)^n}{3^{n-1}}
              \]
              を代入し,分母分子に$3^{n-1}$を掛けて整理すると,
              \[
                a_n=\frac{2\left(3^{n-1}+(-1)^n\right)}{3^{n-1}-(-1)^n}
              \]
              を得る.
        \item $\displaystyle a_1=5,\, a_{n+1}=\dfrac{4a_{n}-9}{a_{n}-2}$
              不動点の方程式は
              \[
                x=\frac{4x-9}{x-2}
              \]
              であるから,その解は3.
              \[
              b_n=a_n-3
              \]
              とおくと
              \begin{align*}
                b_{n+1}&=a_{n+1}-3\\
                &=\dfrac{4a_{n}-9}{a_{n}-2}-3\\
                &=\dfrac{a_{n}-3}{a_{n}-2}\\
                &=\dfrac{b_{n}}{b_{n}+1}
              \end{align*}
              逆数を取って
              \begin{align*}
                \frac{1}{b_{n+1}}&=\dfrac{b_{n}+1}{b_{n}}
              \end{align*}
              $b_1=\dfrac{1}{2}$より,
              \begin{align*}
                &\frac{1}{b_n}=\frac{1}{2}+(n-1)\cdot 1\\
                &\frac{1}{b_n}=\frac{2n-1}{2}\\
                &b_n=\frac{2}{2n-1}\\
                &a_n-3=\frac{2}{2n-1}\\
                &a_n=\frac{6n-1}{2n-1}
              \end{align*}
        \item $\displaystyle a_1=1,\, a_{n+1}=1+\frac{1}{a_n}$\\
              式を変形して
              $\displaystyle a_1=1,\, a_{n+1}=\frac{a_n+1}{a_n}$
              不動点の方程式は
              \[
                x=\frac{x+1}{x}\tag{1}
              \]
              であるから,その解をそれぞれ$\alpha,\beta$とすると
              \[
                \alpha=\frac{1+\sqrt{5}}{2},\ \beta=\frac{1-\sqrt{5}}{2}
              \]
              ここで(1)について,その解$\alpha,\beta$では
              \[
                \alpha+\beta=1,\, \quad \alpha\beta=-1\tag{2}
              \]
              が解と係数の関係から成り立つ.従って
              \begin{align*}
                b_{n+1}&=\frac{a_{n+1}-\alpha}{a_{n+1}-\beta}\\
                &=\cfrac{\cfrac{a_n+1}{a_n}-\alpha}{\cfrac{a_n+1}{a_n}-\beta}\\
                &=\frac{a_n+1-\alpha a_n}{a_n+1-\beta a_n}\\
                &=\frac{(1-\alpha)a_n+1}{(1-\beta)a_n+1}\tag{3}
              \end{align*}
              (2)より,(3)は
              \begin{align*}
                \frac{(1-\alpha)a_n+1}{(1-\beta)a_n+1}&=\frac{\beta a_n+1}{\alpha a_n+1}\\
                &=\frac{\beta a_n-\alpha\beta}{\alpha a_n-\alpha\beta}\\
                &=\frac{\beta (a_n-\alpha)}{\alpha (a_n-\beta)}\\
                &=b_{n+1}=\frac{\beta}{\alpha}b_n\\
              \end{align*}
              初項$b_1$は,
              \begin{align*}
                b_1=\frac{a_1-\alpha}{a_1-\beta}=\frac{1-\alpha}{1-\beta}=\frac{\beta}{\alpha}
              \end{align*}
              従って
              \[
                b_n=\left(\frac{\beta}{\alpha}\right)^n
              \]
              次に
              \[
                b_{n}=\frac{a_{n}-\alpha}{a_{n}-\beta}
              \]
              を$a_n$について解くと,
              \[
                a_{n}=\frac{\beta b_{n}-\alpha}{b_{n}-1}
              \]
              従って
              \begin{align*}
                a_{n}&=\cfrac{\beta \left(\cfrac{\beta}{\alpha}\right)^n-\alpha}{\left(\cfrac{\beta}{\alpha}\right)^n-1}  \\
                &=\frac{\beta^{n+1}-\alpha^{n+1}}{\beta^n-\alpha^n}\\
                &=\cfrac{\left(\cfrac{1-\sqrt{5}}{2}\right)^{n+1}-\left(\cfrac{1+\sqrt{5}}{2}\right)^{n+1}}{\left(\cfrac{1-\sqrt{5}}{2}\right)^n-\left(\cfrac{1+\sqrt{5}}{2}\right)^n}
              \end{align*}        
      \end{enumerate}
    余談にはなるが,最後の問題はビネの公式と呼ばれる,フィボナッチ数列とよく関わりのある概念を使えば簡単に解ける.
    ここで語ってしまうと追加で何ページも必要になってしまうため,割愛する.

    加えて,このタイプの漸化式は帰納法によって示すことが比較的容易なパターンが存在する.
    例えば(3)は解答の一般項から分かるように,分母分子それぞれが一次関数の形で書かれているため,想像しやすい.
    $a_1$から$a_3$ぐらいまでを求めさせてから帰納法で示させるような誘導問題が出た時は帰納法を使って解くことも視野に入れておくといいだろう.
    以下にその例を示す.
    \begin{quote}
      \textbf{【例題】}
      \begin{quote}
        以下の漸化式
        \[
          a_1=5,\, a_{n+1}=\dfrac{4a_{n}-9}{a_{n}-2}
        \]
        を満たす数列$a_n$の一般項を求めよ.
      \end{quote}
      \textbf{【方針】}
      \begin{quote}
        実験的に代入計算をしてみると,
        \[
          a_1=\frac{5}{1},\,a_2=\frac{11}{3},\,a_3=\frac{17}{5},\,a_4=\frac{23}{7},\,a_5=\frac{29}{9}
        \]
        となる.分母が2づつ,分子が6づつ増えていることから,一般項は
        \[
          a_n=\frac{6n-1}{2n-1}
        \]
        と予想する.
      \end{quote}
      \textbf{【解答】}
      \begin{quote}
        一般項$a_n$が
        \[
          a_n=\frac{6n-1}{2n-1}
        \]
        であることを示す.\\
        $n=1$の時,
        \[
          \frac{6-1}{2-1}=5
        \]
        従って成立する.\\
        とある正の整数$k$によって,$n=k$が成り立つ時,$n=k+1$は,
        \begin{align*}
          a_{k+1}&=\frac{4a_{k}-9}{a_{k}-2}\\
          &=\cfrac{4\cdot \cfrac{6k-1}{2k-1}-9}{\cfrac{6k-1}{2k-1}-2}\\
          &=\frac{4(6k-1)-9(2k-1)}{6k-1-2(2k-1)}\\
          &=\frac{6k+5}{2k+1}\\
          &=\frac{6(k+1)-1}{2(k+1)-1}
        \end{align*}
        従って成立する.
        以上より
        \[
          a_n=\frac{6n-1}{2n-1}
        \]
        が示された.\qed
      \end{quote}


    \end{quote}

  \subsection{隣接3項間漸化式型 \, $\displaystyle a_{n+2}+pa_{n+1}+qa_n=0$}
    \subsubsection{二次特性方程式型 \, $\displaystyle a_{n+2}+pa_{n+1}+qa_n=0$}
      節は変われど,特性方程式を使って考える.
      この小節でも0除算が発生しうる状況があるが,そのような場合はそもそもそのような漸化式が成立しないか,別の解法で解くものとなっているため数学的な問題はない.
      結局,大方の場合,等比数列に変形できればいいのだから,これを等比数列に変形できないかといくつかの試行錯誤を重ねるのが王道の解き方である.
      結果から言ってしまえば,特性方程式型のところで少し触れた式を生み出すことで等比数列を作ることができる.
      \[
        a_{n+2}+pa_{n+1}+qa_n=0
      \]
      について,
      \[
        a_{n+2}+\alpha a_{n+1}=\beta(a_{n+1}+\alpha a_{n})
      \]
      となるような$\alpha,\beta$の組を探せば良い.
      変形すると
      \begin{align*}
        &a_{n+2}+\alpha a_{n+1}=\beta(a_{n+1}+\alpha a_{n})\\
        &\therefore \quad a_{n+2}+(\alpha-\beta)a_{n+1}-\alpha\beta a_n=0
      \end{align*}
      である.ここで$\beta'=-\beta$と置き,不動点を考えているから$a_n$に関する項を$x$と置けば,
      \[
        x^2+(\alpha+\beta')x+\alpha\beta'=0      
      \]
      となり,まさに二次方程式の解を求めることに一致する.
      従って形式的には
      \[
        a_{n+2}+pa_{n+1}+qa_n=0
      \]
      については,
      \[
        x^2+px+q=0
      \]
      について解くことで既知の形に帰着する.

      \,$x^2+px+q=0$の2解を$\alpha,\beta$とすると,元の漸化式は
      \[
        \begin{cases}
          a_{n+1}-\alpha a_n=\beta(a_n-\alpha a_{n-1})\\
          a_{n+1}-\beta a_n=\alpha(a_n-\beta a_{n-1})
        \end{cases}
      \]
      という2つの等比数列を生む(添字は1つずれているが本質的には先程と同じ議論である).
      すなわち,\,$b_n=a_{n+1}-\alpha a_n$は公比$\beta$の等比数列,\,$c_n=a_{n+1}-\beta a_n$は公比$\alpha$の等比数列になる.
      $\alpha\ne\beta$であれば,この2つの式から$a_{n+1}$を消去して$a_n$について解けばよい.
      また,\,$\alpha=\beta$(重解)の場合は等比数列が1つしか作れないため,\,$b_n=a_{n+1}-\alpha a_n$が公比$\alpha$の等比数列であることを使い,
      両辺を$\alpha^{n+1}$で割って階差型に帰着させることになる.

      \noindent \textbf{【例題】}
      \begin{enumerate}[label=(\arabic*),parsep=3pt]
        \item $\displaystyle a_1=1,\,a_2=1,\, a_{n+2}=5a_{n+1}-6a_n$
        \item $\displaystyle a_1=1,\,a_2=4,\, a_{n+2}=4a_{n+1}-4a_n$
        \item $\displaystyle a_1=2,\,a_2=1,\, a_{n+2}=a_{n+1}+6a_n$
      \end{enumerate}

      \noindent\textbf{【方針】}
      \begin{enumerate}[label=(\arabic*),parsep=3pt]
        \item 特性方程式$x^2-5x+6=0$の解は異なる2つの実数解です.素直に2本の等比数列を作りましょう.
        \item 特性方程式$x^2-4x+4=0$を解くと重解になります.等比数列は1本しか作れないことに注意しましょう.
        \item 特性方程式の解の1つが負になるパターンです.符号に注意して計算しましょう.
      \end{enumerate}

      \noindent\textbf{【解答】}
      \begin{enumerate}[label=(\arabic*),parsep=3pt]
        \item $\displaystyle a_n=2^n-3^{n-1}$
        \item $\displaystyle a_n=n\cdot2^{n-1}$
        \item $\displaystyle a_n=(-2)^{n-1}+3^{n-1}$
      \end{enumerate}

      \noindent\textbf{【解法】}
      \begin{enumerate}[label=(\arabic*),parsep=3pt]
        \item $\displaystyle a_1=1,\,a_2=1,\, a_{n+2}=5a_{n+1}-6a_n$\\
              特性方程式$x^2-5x+6=0$を解くと,\\
              $(x-2)(x-3)=0$より$\alpha=2,\beta=3$である.
              \begin{align*}
                &a_{n+2}-2a_{n+1}=3(a_{n+1}-2a_n)\\
                &a_{n+2}-3a_{n+1}=2(a_{n+1}-3a_n)
              \end{align*}
              なので,\,$b_n=a_{n+1}-2a_n,\ c_n=a_{n+1}-3a_n$と置くと,
              \begin{align*}
                &b_n=b_1\cdot3^{n-1},\quad b_1=a_2-2a_1=-1\\
                &c_n=c_1\cdot2^{n-1},\quad c_1=a_2-3a_1=-2\\
                &\therefore \quad b_n=-3^{n-1},\quad c_n=-2^n
              \end{align*}
              $b_n-c_n=(a_{n+1}-2a_n)-(a_{n+1}-3a_n)=a_n$であるから,
              \[
                a_n=b_n-c_n=-3^{n-1}-(-2^n)=2^n-3^{n-1}
              \]
        \item $\displaystyle a_1=1,\,a_2=4,\, a_{n+2}=4a_{n+1}-4a_n$\\
              特性方程式$x^2-4x+4=0$を解くと,\,$(x-2)^2=0$より重解$\alpha=2$である.
              \,$b_n=a_{n+1}-2a_n$と置くと,
              \begin{align*}
                &b_{n+1}=a_{n+2}-2a_{n+1}\\
                &=(4a_{n+1}-4a_n)-2a_{n+1}\\
                &\qquad=2(a_{n+1}-2a_n)=2b_n\\
                &b_1=a_2-2a_1=2\\
                &\therefore \quad b_n=2\cdot2^{n-1}=2^n
              \end{align*}
              したがって$a_{n+1}-2a_n=2^n$である.両辺を$2^{n+1}$で割ると,
              \begin{align*}
                &\frac{a_{n+1}}{2^{n+1}}-\frac{a_n}{2^n}=\frac12\\
                &\text{$c_n=\dfrac{a_n}{2^n}$と置くと,}\\
                &c_{n+1}=c_n+\frac12,\quad c_1=\frac{a_1}{2}=\frac12\\
                &\therefore \quad c_n=\frac12+(n-1)\cdot\frac12=\frac n2\\
                &\therefore \quad a_n=2^n\cdot c_n=n\cdot2^{n-1}
              \end{align*}
        \item $\displaystyle a_1=2,\,a_2=1,\, a_{n+2}=a_{n+1}+6a_n$\\
              特性方程式$x^2-x-6=0$を解くと,\,$(x-3)(x+2)=0$より$\alpha=3,\beta=-2$である.
              \begin{align*}
                &a_{n+2}-3a_{n+1}=-2(a_{n+1}-3a_n)\\
                &a_{n+2}+2a_{n+1}=3(a_{n+1}+2a_n)
              \end{align*}
              なので,\,$b_n=a_{n+1}-3a_n,\ c_n=a_{n+1}+2a_n$と置くと,
              \begin{align*}
                &b_n=b_1\cdot(-2)^{n-1},\quad b_1=a_2-3a_1=-5\\
                &c_n=c_1\cdot3^{n-1},\quad c_1=a_2+2a_1=5\\
                &\therefore \quad b_n=-5(-2)^{n-1},\quad c_n=5\cdot3^{n-1}
              \end{align*}
              $c_n-b_n=(a_{n+1}+2a_n)-(a_{n+1}-3a_n)=5a_n$であるから,
              \begin{align*}
                5a_n&=5\cdot3^{n-1}-\{-5(-2)^{n-1}\}\\
                &=5\cdot3^{n-1}+5(-2)^{n-1}\\
                \therefore\quad a_n&=(-2)^{n-1}+3^{n-1}
              \end{align*}
      \end{enumerate}

    \subsubsection{等差二次特性方程式型 \, $\displaystyle a_{n+2}+pa_{n+1}+qa_n=r$}
      今度は右辺が0ではなく$r$となった.
      つまり$r$だけズレたのである.
      ズレたらもとに戻す,つまり不動点を探すのである.
      \[
        b_n=a_n-\alpha
      \]
      となる$\alpha$を求めると,
      \begin{align*}
        &(a_{n+2}-\alpha)+p(a_{n+1}-\alpha)+q(a_n-\alpha)=0\\
        &a_{n+2}+pa_{n+1}+qa_n=\alpha+p\alpha+q\alpha\\
        &\therefore \quad r=\alpha +p\alpha +q\alpha\\
        &\Leftrightarrow \quad \alpha=\frac{r}{p+q+1}
      \end{align*}
      あとは前節で述べた隣接3項間漸化式の問題である.

      \noindent \textbf{【例題】}
      \begin{enumerate}[label=(\arabic*),parsep=3pt]
        \item $\displaystyle a_1=2,\,a_2=2,\, a_{n+2}=5a_{n+1}-6a_n+2$
        \item $\displaystyle a_1=1,\,a_2=2,\, a_{n+2}=4a_{n+1}-4a_n+3$
      \end{enumerate}

      \noindent\textbf{【方針】}
      \begin{enumerate}[label=(\arabic*),parsep=3pt]
        \item まず不動点$\alpha$を求め,\,$b_n=a_n-\alpha$と置いて前節の二次特性方程式型に帰着させましょう.
        \item こちらは重解のパターンです.同じ手順で$\alpha$を求めてから前節の重解の場合の処理を行いましょう.
      \end{enumerate}

      \noindent\textbf{【解答】}
      \begin{enumerate}[label=(\arabic*),parsep=3pt]
        \item $\displaystyle a_n=1+2^n-3^{n-1}$
        \item $\displaystyle a_n=3+2^{n-2}(3n-7)$
      \end{enumerate}

      \noindent\textbf{【解法】}
      \begin{enumerate}[label=(\arabic*),parsep=3pt]
        \item $\displaystyle a_1=2,\,a_2=2,\, a_{n+2}=5a_{n+1}-6a_n+2$
              \[
                a_{n+2}=5a_{n+1}-6a_n+2 \iff (a_{n+2}-\alpha)-5(a_{n+1}-\alpha)+6(a_n-\alpha)=0 
              \]
              となる$\alpha$を求めると,
              \[
                \alpha=1
              \]
              $b_n=a_n-1$と置くと,
              \begin{align*}
                &b_{n+2}=5b_{n+1}-6b_n\\
                &b_1=1,\, b_2=1
              \end{align*}
              （答案中略。前節の例題(1)と同じ漸化式に帰着する.全く同じ計算を$b_n$に対して行う)
              \[
                b_n=2^n-3^{n-1}
              \]
              を得るから,
              \[
                a_n=1+b_n=1+2^n-3^{n-1}
              \]
        \item $\displaystyle a_1=1,\,a_2=2,\, a_{n+2}=4a_{n+1}-4a_n+3$
              \[
                a_{n+2}=4a_{n+1}-4a_n+3 \iff (a_{n+2}-\alpha)-4(a_{n+1}-\alpha)+4(a_n-\alpha)=0 
              \]
              となる$\alpha$を求めると,
              \[
                \alpha=3
              \]
              なので,\,$b_n=a_n-3$と置くと,
              \begin{align*}
                &b_{n+2}=4b_{n+1}-4b_n\\
                &b_1=a_1-3=-2\\
                &b_2=a_2-3=-1                
              \end{align*}
              特性方程式$x^2-4x+4=0$は重解$\alpha'=2$を持つから,\,$c_n=b_{n+1}-2b_n$と置くと,
              \begin{align*}
                &c_{n+1}=2c_n\\
                &c_1=b_2-2b_1=-1-2(-2)=3\\
                &\therefore \quad c_n=3\cdot2^{n-1}
              \end{align*}
              したがって$b_{n+1}-2b_n=3\cdot2^{n-1}$である.両辺を$2^{n+1}$で割ると,
              \begin{align*}
                &\frac{b_{n+1}}{2^{n+1}}-\frac{b_n}{2^n}=\frac34\\
                &\text{$f_n=\dfrac{b_n}{2^n}$と置くと,}\\
                &f_{n+1}=f_n+\frac34,\quad f_1=\frac{b_1}{2}=-1\\
                &\therefore \quad f_n=-1+(n-1)\cdot\frac34=\frac{3n-7}{4}\\
                &\therefore \quad b_n=2^n\cdot f_n=2^{n-2}(3n-7)
              \end{align*}
              よって,
              \[
                a_n=3+b_n=3+2^{n-2}(3n-7)
              \]
      \end{enumerate}
  \subsection{総和型}
    \subsubsection{総和と一般項の漸化式型 \, $S_n=pa_n+f(n)$}
      ある漸化式によって定められる数列$a_n$によって,
      \[
        S_n=\sum_{k=1}^{n} a_k
      \]
      として定められる数列$S_n$について考えよう.
      $S_n$と$a_n$の関係を表した
      \[
        S_n=pa_n+f(n)
      \]
      のような漸化式を考える.
      このような漸化式は今まで取り扱った形にはない,未知のものである.
      であれば既知の形にするのが数学である.
      これまでのことをきちんと理解していれば,$a_{n+1}$と$a_n$の関係が分かっている時には解くことができる.
      $S_n$をうまく変形して$a_{n}$関連だけの式を産むには,\,$p \neq 1$の時,
      \begin{align*}
        &S_n=pa_n+f(n) \tag{1}\\
        &S_{n+1}=pa_{n+1}+f(n+1) \tag{2}
      \end{align*}
      $(2)-(1)$をして
      \begin{align*}
        &a_{n+1}=pa_{n+1}-pa_n+f(n+1)-f(n)\\
        &a_{n+1}-pa_{n+1}=-pa_n+f(n+1)-f(n)\\
        &a_{n+1}(1-p)=-pa_n+f(n+1)-f(n)\\
        &a_{n+1}=\frac{-pa_n+f(n+1)-f(n)}{(1-p)}\\
        &a_{n+1}=\frac{p}{p-1}a_n+\frac{1}{1-p}\{f(n+1)-f(n)\}
      \end{align*}
      $\displaystyle g(n)=\frac{1}{1-p}\{f(n+1)-f(n)\}$とおくと,
      \begin{align*}
        &a_{n+1}=\frac{p}{p-1}a_n+g(n)
      \end{align*}
      従って漸化式は階差数列型に帰着する.$p=1$のときも同様の流れで求めることができる.

      \noindent \textbf{【例題】}
      \begin{enumerate}[label=(\arabic*),parsep=3pt]
        \item $\displaystyle S_n=2a_n-n$
        \item $\displaystyle S_n=\frac12a_n+n$
      \end{enumerate}

      \noindent\textbf{【方針】}
      \begin{enumerate}[label=(\arabic*),parsep=3pt]
        \item まず$n=1$を代入して$a_1$を求めてから,本文の変形をなぞって$a_{n+1}$と$a_n$だけの漸化式を作りましょう.
        \item 同様です.できた漸化式は特性方程式型に帰着します.
      \end{enumerate}

      \noindent\textbf{【解答】}
      \begin{enumerate}[label=(\arabic*),parsep=3pt]
        \item $\displaystyle a_n=2^n-1$
        \item $\displaystyle a_n=1+(-1)^{n-1}$
      \end{enumerate}

      \noindent\textbf{【解法】}
      \begin{enumerate}[label=(\arabic*),parsep=3pt]
      \item $\displaystyle S_n=2a_n-n$

          $n=1$を代入すると,
          \[
            S_1=a_1=2a_1-1
          \]
          より
          \[
            a_1=1.
          \]
          また,
          \begin{align*}
            a_{n+1}
              &=S_{n+1}-S_n\\
              &=\{2a_{n+1}-(n+1)\}-(2a_n-n)\\
              &=2a_{n+1}-2a_n-1.
          \end{align*}
          したがって,
          \[
            a_{n+1}=2a_n+1.
          \]
          ここで $b_n=a_n+1$ とおくと,
          \begin{align*}
            b_{n+1}
              &=a_{n+1}+1\\
              &=2a_n+2\\
              &=2(a_n+1)\\
              &=2b_n.
          \end{align*}
          また,
          \[
            b_1=a_1+1=2
          \]
          であるから,
          \[
            b_n=2^n.
          \]
          よって,
          \[
            a_n=2^n-1.
          \]
      \item $\displaystyle S_n=\frac12a_n+n$
          $n=1$を代入すると,
          \[
            S_1=a_1=\frac12a_1+1
          \]
          より
          \[
            a_1=2.
          \]
          \begin{align*}
            a_{n+1}
              &=S_{n+1}-S_n\\
              &=\left\{\frac12a_{n+1}+(n+1)\right\}
                -\left(\frac12a_n+n\right)\\
              &=\frac12a_{n+1}-\frac12a_n+1.
          \end{align*}
          したがって,
          \[
            a_{n+1}=-a_n+2.
          \]

          ここで $b_n=a_n-1$ とおくと,
          \begin{align*}
            b_{n+1}
              &=a_{n+1}-1\\
              &=-a_n+1\\
              &=-(a_n-1)\\
              &=-b_n.
          \end{align*}
          また,
          \[
            b_1=a_1-1=1
          \]
          であるから,
          \[
            b_n=(-1)^{n-1}.
          \]
          よって,
          \[
            a_n=1+(-1)^{n-1}.
          \]

      \end{enumerate}

    \subsubsection{純総和型 \, $S_{n+1}=pS_n+f(n) $}
      両辺に総和がきてもやることは一緒で,1つずらした総和と元の式を足し引きしてやることで,\,$a_n$や$a_{n+1}$の項を作れば良い.
      \begin{align*}
        &S_{n+1}=pS_n+f(n) \tag{1}\\
        &S_{n}=pS_{n-1}+f(n-1) \tag{2}\\
        &\text{$(1)-(2)$をして,}\\
        &a_{n+1}=pa_{n}+f(n)-f(n-1)
      \end{align*}
      あとはただの階差数列型の計算である.

      \noindent \textbf{【例題】}
      \begin{enumerate}[label=(\arabic*),parsep=3pt]
        \item $\displaystyle S_1=1,\ S_{n+1}=2S_n+n$
      \end{enumerate}

      \noindent\textbf{【方針】}
      \begin{enumerate}[label=(\arabic*),parsep=3pt]
        \item 本文の式は$n\ge2$でしか使えないことに注意し,\,$a_2$だけは$S_2=2S_1+f(1)$から別に求めましょう.
      \end{enumerate}

      \noindent\textbf{【解答】}
      \begin{enumerate}[label=(\arabic*),parsep=3pt]
        \item $\displaystyle a_1=1,\quad a_n=3\cdot2^{n-2}-1\ (n\ge2)$
      \end{enumerate}

      \noindent\textbf{【解法】}
      \begin{enumerate}[label=(\arabic*),parsep=3pt]
      \item $\displaystyle S_1=1,\ S_{n+1}=2S_n+n$
          $a_1=S_1=1$である.
          また,
          \[
            S_{n+1}=a_1+a_2+\cdots+a_n+a_{n+1},
            \qquad
            S_n=a_1+a_2+\cdots+a_n
          \]
          であるから,
          \[
            S_{n+1}-S_n=a_{n+1}.
          \]
          したがって,
          \begin{align*}
            a_{n+1}
              &=S_{n+1}-S_n\\
              &=(2S_n+n)-S_n\\
              &=S_n+n.
          \end{align*}
          ここで,
          \[
            S_n=2S_{n-1}+(n-1)
          \]
          より,
          \begin{align*}
            a_{n+1}
              &=2S_{n-1}+(n-1)+n\\
              &=2S_{n-1}+2n-1.
          \end{align*}
          一方,
          \[
            a_n=S_n-S_{n-1}
          \]
          であるから,
          \[
            S_n=2S_{n-1}+(n-1)
          \]
          を用いて,
          \[
            a_n=S_{n-1}+n-1,
          \]
          したがって
          \[
            S_{n-1}=a_n-(n-1).
          \]
          これを上式に代入すると,
          \begin{align*}
            a_{n+1}
              &=2\{a_n-(n-1)\}+2n-1\\
              &=2a_n+1.
          \end{align*}

          なお,$n=1$のときは
          \[
            a_2=S_2-S_1=(2S_1+1)-S_1=2
          \]
          であり,上の漸化式
          \[
            a_{n+1}=2a_n+1
          \]
          は$n\ge1$で成り立つ.

          ここで $b_n=a_n+1$ とおくと,
          \begin{align*}
            b_{n+1}
              &=a_{n+1}+1\\
              &=2a_n+2\\
              &=2(a_n+1)\\
              &=2b_n.
          \end{align*}
          また,
          \[
            b_1=a_1+1=2
          \]
          であるから,
          \[
            b_n=2^n.
          \]
          よって,
          \[
            a_n=2^n-1.
          \]
      \end{enumerate}

  \subsection{連立型}
    \subsubsection{加減法型}
      一般形が
      \[
        \begin{cases}
          a_{n+1}=pa_{n}+qb_{n} \\
          b_{n+1}=qa_{n}+pb_{n}
        \end{cases}
      \]
      の形で表される漸化式の解き方について考えよう.
      といってもここまで読んだ諸君らにとっては,直感的に与えられた式を足し引きすればいいことが分かるだろう.
      実際にすすめていくと,
      \begin{align*}
        &a_{n+1}+b_{n+1}=(p+q)(a_{n}+b_{n})\\
        &a_n+b_n=(a_1+b_1)(p+q)^{n-1}\tag{1}
      \end{align*}
      また,
      \begin{align*}
        &a_{n+1}-b_{n+1}=(p-q)(a_{n}-b_{n})\\
        &a_n-b_n=(a_1-b_1)(p-q)^{n-1}\tag{2}
      \end{align*}
      $(1)+(2)$をすると
      \begin{align*}
        2a_n=(a_1+b_1)(p+q)^{n-1}+(a_1-b_1)(p-q)^{n-1}\\
        \therefore \quad a_n=\frac{(a_1+b_1)(p+q)^{n-1}+(a_1-b_1)(p-q)^{n-1}}{2}\\
      \end{align*}
      $(1)-(2)$をすると
      \begin{align*}
        2b_n=(a_1+b_1)(p+q)^{n-1}-(a_1-b_1)(p-q)^{n-1}\\
        \therefore \quad b_n=\frac{(a_1+b_1)(p+q)^{n-1}-(a_1-b_1)(p-q)^{n-1}}{2}
      \end{align*}
      別解及び演習は一般型の方で行う.
    \subsubsection{一般型}
      連立一般型の漸化式は,
      \[
        \begin{cases}
          a_{n+1}=pa_{n}+qb_{n} \\
          b_{n+1}=ra_{n}+sb_{n}
        \end{cases}
      \]
      で表される漸化式である.
      いつものように等比数列型への帰着を目指そう.
      $a_n$と$b_n$が互いに影響し合っているため,独立させて解くことは難しそうだ.
      ならばそれぞれについて求めるのではなく,いっそのこと合わせて求めてしまおうというのが解法の思想である.
      その流れは,
      \[
        \begin{array}{rrl}
          & a_{n+1}={} &pa_{n}+qb_{n}\\
          +\,) & b_{n+1}={} & ra_{n}+sb_{n}\\ \hline
          & a_{n+1}+b_{n+1}={} & (p+r)a_{n}+(q+s)b_{n}
        \end{array}
      \]
      について
      \[
        \alpha a_{n+1}+\beta b_{n+1}
        =\gamma(\alpha a_n+\beta b_n)
      \]
      を満たすようにしてやればよい.
      実際に左辺を計算すると,
      \begin{align*}
        \alpha a_{n+1}+\beta b_{n+1}&=\alpha(pa_n+qb_n)+\beta(ra_n+sb_n)\\
        &=(p\alpha+r\beta)a_n+(q\alpha+s\beta)b_n.
      \end{align*}
      一方,
      \[
        \gamma(\alpha a_n+\beta b_n)=\gamma\alpha a_n+\gamma\beta b_n
      \]
      である.
      したがって,
      \[
      \begin{cases}
        p\alpha+r\beta=\gamma\alpha\\
        q\alpha+s\beta=\gamma\beta
      \end{cases}
      \]
      となるように$\alpha,\beta,\gamma$を決めればよい.

      ここで,いきなり3つの文字をすべて求めようとすると少し複雑に見える.
      そこで,まず$\gamma$について考える.

      上の2式を変形すると,
      \[
      \begin{cases}
        (p-\gamma)\alpha+r\beta=0\\
        q\alpha+(s-\gamma)\beta=0
      \end{cases}
      \]
      である.
      $\alpha,\beta$は同時に$0$ではないので,この連立方程式が$\alpha,\beta$について解を持つためには,
      \[
        (p-\gamma)(s-\gamma)-qr=0
      \]
      でなければならない.
      整理すると,
      \[
        \gamma^2-(p+s)\gamma+(ps-qr)=0
      \]
      となる.
      この二次方程式を,この連立漸化式の特性方程式と呼ぶ.
      特性方程式の解を$\gamma_1,\gamma_2$とする.
      まず$\gamma=\gamma_1$を上の連立方程式に代入し,
      \[
        \begin{cases}
          (p-\gamma_1)\alpha+r\beta=0\\
          q\alpha+(s-\gamma_1)\beta=0
        \end{cases}
      \]
      を満たす$\alpha,\beta$を1組求める.
      すると,
      \[
        c_n=\alpha a_n+\beta b_n
      \]
      と置くことで,
      \begin{align*}
        c_{n+1}&=\alpha a_{n+1}+\beta b_{n+1}\\
        &=\gamma_1(\alpha a_n+\beta b_n)\\
        &=\gamma_1c_n
      \end{align*}
      となる.
      したがって,
      \[
        c_n=c_1\gamma_1^{\,n-1}
      \]
      という等比数列が得られる.
      同様に,$\gamma=\gamma_2$に対しても$\alpha,\beta$を求めれば,
      \[
        d_n=\alpha' a_n+\beta' b_n
      \]
      について
      \[
        d_n=d_1\gamma_2^{\,n-1}
      \]
      という等比数列が得られる.
      なお,ここで得られた2つの等比数列から$a_n,b_n$を求めるには,
      \[
        \begin{cases}
        c_n=\alpha a_n+\beta b_n\\
        d_n=\alpha'a_n+\beta'b_n
        \end{cases}
      \]
      を$a_n,b_n$について解けばよい.
      このように,$a_n,b_n$をそれぞれ直接求めるのではなく,$a_n,b_n$の和の形を作ることで,連立していた漸化式を等比数列へと分離することができる.

      \noindent \textbf{【例題】}
      \begin{enumerate}[label=(\arabic*),parsep=3pt]
        \item $\displaystyle a_1=1,\,b_1=1.\ \begin{cases}a_{n+1}=3a_n+b_n\\ b_{n+1}=a_n+3b_n\end{cases}$
        \item $\displaystyle a_1=1,\,b_1=0,\ \begin{cases}a_{n+1}=3a_n+b_n\\ b_{n+1}=2a_n+2b_n\end{cases}$
      \end{enumerate}

      \noindent\textbf{【方針】}
      \begin{enumerate}[label=(\arabic*),parsep=3pt]
        \item 前節で述べた問題です.一般型,加減法型,どちらのやり方でも問題ありません.
        \item これはこの節で述べたものです.
      \end{enumerate}

      \noindent\textbf{【解答】}
      \begin{enumerate}[label=(\arabic*),parsep=3pt]
        \item $\displaystyle a_n=\frac{2\cdot4^{n-1}+1}{3},\qquad b_n=\frac{2\cdot4^{n-1}-2}{3}$
        \item $\displaystyle a_n=4^{n-1},\qquad b_n=a_n=4^{n-1}$
      \end{enumerate}

      \noindent\textbf{【解法】}
      \begin{enumerate}[label=(\arabic*),parsep=3pt]
      \item $\displaystyle a_1=1,\,b_1=1.\ \begin{cases}a_{n+1}=3a_n+b_n\\ b_{n+1}=a_n+3b_n\end{cases}$\\
            \begin{enumerate}
              \item 【解法1】\\
                    $\alpha a_{n+1}+\beta b_{n+1}$を計算すると,
                    \begin{align*}
                      \alpha a_{n+1}+\beta b_{n+1}
                        &=\alpha(3a_n+b_n)+\beta(a_n+3b_n)\\
                        &=(3\alpha+\beta)a_n+(\alpha+3\beta)b_n.
                    \end{align*}
                    これが
                    \[
                      \gamma(\alpha a_n+\beta b_n)
                      =\gamma\alpha a_n+\gamma\beta b_n
                    \]
                    と一致するためには,
                    \[
                    \begin{cases}
                      (3-\gamma)\alpha+\beta=0\\
                      \alpha+(3-\gamma)\beta=0
                    \end{cases}
                    \]
                    でなければならない.

                    $\alpha,\beta$が同時に$0$ではないため,
                    \[
                      (3-\gamma)^2-1=0
                    \]
                    より,
                    \[
                      \gamma^2-6\gamma+8=0,
                      \qquad
                      (\gamma-2)(\gamma-4)=0.
                    \]
                    したがって,
                    \[
                      \gamma_1=2,\qquad\gamma_2=4
                    \]
                    である.

                    まず$\gamma=2$とすると,
                    \[
                    \begin{cases}
                      \alpha+\beta=0\\
                      \alpha+\beta=0
                    \end{cases}
                    \]
                    となるので,
                    \[
                      \alpha=1,\qquad\beta=-1
                    \]
                    と取ることができる.

                    そこで,
                    \[
                      c_n=a_n-b_n
                    \]
                    とおくと,
                    \begin{align*}
                      c_{n+1}
                        &=a_{n+1}-b_{n+1}\\
                        &=(3a_n+b_n)-(a_n+3b_n)\\
                        &=2(a_n-b_n)\\
                        &=2c_n.
                    \end{align*}
                    また,
                    \[
                      c_1=a_1-b_1=0
                    \]
                    であるから,
                    \[
                      c_n=0.
                    \]
                    したがって,
                    \[
                      a_n-b_n=0.
                    \]

                    次に$\gamma=4$とすると,
                    \[
                    \begin{cases}
                      -\alpha+\beta=0\\
                      \alpha-\beta=0
                    \end{cases}
                    \]
                    となるので,
                    \[
                      \alpha=1,\qquad\beta=1
                    \]
                    と取ることができる.

                    そこで,
                    \[
                      d_n=a_n+b_n
                    \]
                    とおくと,
                    \begin{align*}
                      d_{n+1}
                        &=a_{n+1}+b_{n+1}\\
                        &=(3a_n+b_n)+(a_n+3b_n)\\
                        &=4(a_n+b_n)\\
                        &=4d_n.
                    \end{align*}
                    また,
                    \[
                      d_1=a_1+b_1=2
                    \]
                    であるから,
                    \[
                      d_n=2\cdot4^{n-1}.
                    \]

                    以上より,
                    \[
                    \begin{cases}
                      a_n-b_n=0\\
                      a_n+b_n=2\cdot4^{n-1}
                    \end{cases}
                    \]
                    となる.

                    2式を辺々加えると,
                    \[
                      2a_n=2\cdot4^{n-1}
                    \]
                    であるから,
                    \[
                      a_n=4^{n-1}.
                    \]
                    また,
                    \[
                      b_n=a_n=4^{n-1}.
                    \]
              \item 【解法2】\\
                    2つの式を辺々加えると,
                    \begin{align*}
                      a_{n+1}+b_{n+1}
                        &=(3a_n+b_n)+(a_n+3b_n)\\
                        &=4(a_n+b_n).
                    \end{align*}
                    よって,
                    \[
                      a_n+b_n=(a_1+b_1)4^{n-1}=2\cdot4^{n-1}. \tag{1}
                    \]

                    また,2つの式を辺々引くと,
                    \begin{align*}
                      a_{n+1}-b_{n+1}
                        &=(3a_n+b_n)-(a_n+3b_n)\\
                        &=2(a_n-b_n).
                    \end{align*}
                    よって,
                    \[
                      a_n-b_n=(a_1-b_1)2^{n-1}=0. \tag{2}
                    \]

                    \((1)+(2)\)をすると,
                    \[
                      2a_n=2\cdot4^{n-1}
                    \]
                    なので,
                    \[
                      a_n=4^{n-1}.
                    \]

                    また,\((1)-(2)\)をすると,
                    \[
                      2b_n=2\cdot4^{n-1}
                    \]
                    なので,
                    \[
                      b_n=4^{n-1}.
                    \]
            \end{enumerate}
      \item $\displaystyle a_1=1,\,b_1=0,\ \begin{cases}a_{n+1}=3a_n+b_n\\ b_{n+1}=2a_n+2b_n\end{cases}$\\
            $\alpha a_{n+1}+\beta b_{n+1}$を計算すると,
            \begin{align*}
              \alpha a_{n+1}+\beta b_{n+1}
                &=\alpha(3a_n+b_n)+\beta(2a_n+2b_n)\\
                &=(3\alpha+2\beta)a_n+(\alpha+2\beta)b_n.
            \end{align*}
            これが
            \[
              \gamma(\alpha a_n+\beta b_n)
              =\gamma\alpha a_n+\gamma\beta b_n
            \]
            と一致するためには,
            \[
            \begin{cases}
              (3-\gamma)\alpha+2\beta=0\\
              \alpha+(2-\gamma)\beta=0
            \end{cases}
            \]
            でなければならない.

            $\alpha,\beta$が同時に$0$ではないため,
            \[
              (3-\gamma)(2-\gamma)-2=0
            \]
            より,
            \[
              \gamma^2-5\gamma+4=0,
              \qquad
              (\gamma-1)(\gamma-4)=0.
            \]
            したがって,
            \[
              \gamma_1=1,\qquad\gamma_2=4
            \]
            である.

            まず$\gamma=1$とすると,
            \[
            \begin{cases}
              2\alpha+2\beta=0\\
              \alpha+\beta=0
            \end{cases}
            \]
            となるので,
            \[
              \alpha=1,\qquad\beta=-1
            \]
            と取ることができる.

            そこで,
            \[
              c_n=a_n-b_n
            \]
            とおくと,
            \begin{align*}
              c_{n+1}
                &=a_{n+1}-b_{n+1}\\
                &=(3a_n+b_n)-(2a_n+2b_n)\\
                &=a_n-b_n\\
                &=c_n.
            \end{align*}
            また,
            \[
              c_1=a_1-b_1=1
            \]
            であるから,
            \[
              c_n=1.
            \]
            したがって,
            \[
              a_n-b_n=1.
            \]

            次に$\gamma=4$とすると,
            \[
            \begin{cases}
              -\alpha+2\beta=0\\
              \alpha-2\beta=0
            \end{cases}
            \]
            となるので,
            \[
              \alpha=2,\qquad\beta=1
            \]
            と取ることができる.

            そこで,
            \[
              d_n=2a_n+b_n
            \]
            とおくと,
            \begin{align*}
              d_{n+1}
                &=2a_{n+1}+b_{n+1}\\
                &=2(3a_n+b_n)+(2a_n+2b_n)\\
                &=8a_n+4b_n\\
                &=4(2a_n+b_n)\\
                &=4d_n.
            \end{align*}
            また,
            \[
              d_1=2a_1+b_1=2
            \]
            であるから,
            \[
              d_n=2\cdot4^{n-1}.
            \]

            以上より,
            \[
            \begin{cases}
              a_n-b_n=1\\
              2a_n+b_n=2\cdot4^{n-1}
            \end{cases}
            \]
            となる.

            2式を辺々加えると,
            \[
              3a_n=1+2\cdot4^{n-1}
            \]
            であるから,
            \[
              a_n=\frac{2\cdot4^{n-1}+1}{3}.
            \]
            また,
            \[
              b_n=a_n-1
            \]
            より,
            \[
              b_n=\frac{2\cdot4^{n-1}-2}{3}.
            \]
      \end{enumerate}
\chapter{実際の解き方}
  \section{手法まとめ}
  これまで扱った方法をまとめよう.

  \noindent \textbf{基本4パターン}
  \begin{enumerate}
    \item $a_{n+1}=a_n+d$\\
          等差型.\,$d$が$n-1$回足されるから一般項は$a_n=a_1+(n-1)d$.
    \item $a_{n+1}=ra_n$\\
          等比型.\,$r$が$n-1$回掛けられるから一般項は$a_n=a_1\cdot r^{n-1}$.
    \item $a_{n+1}=a_n+f(n)$\\
          階差型.\,$f(n)$が$n$を変化させつつ$n-1$回足されるから一般項は$\displaystyle a_n=a_1+\sum_{k=1}^{n-1} f(k)$.
    \item $a_{n+1}=f(n)a_n$\\
          階比型.\,$f(n)$が$n$を変化させつつ$n-1$回掛けられる一般項は$\displaystyle a_n=a_1\cdot\prod_{k=1}^{n-1} f(k)$.\\
          特に$f(n)=n$や$f(n)=n(n+1)$の場合は一番大きな項の階乗などで割る.分数の場合は掛ける.
  \end{enumerate}

  \noindent \textbf{応用7パターン}
  \begin{enumerate}
    \item $a_{n+1}=pa_n+q$\\
          特性方程式型.$a_{n+1}-\alpha=p(a_n-\alpha)$を満たすような$\alpha$を見つけて等比型へ帰着.（$\alpha=p\alpha+q$を解いて$\alpha$を導出しても良い.）
    \item $a_{n+1}=pa_n+f(n)$\\
          \begin{enumerate}
            \item $a_{n+1}=pa_n+f(n)$
                  階差等比複合型.$a_{n+1}-g(n+1)=p(a_n-g(n))$を満たすような$g(n)$を見つけて等比型へ帰着.
            \item $a_{n+1}=pa_n+q^n$\\
                  $q^n$で割って特性方程式型へ帰着.
          \end{enumerate}
    \item $a_{n+1}=pa_n^q$\\
          対数型.$p$や$p$に関連する数を底とする対数を取って特性方程式型へ帰着.
    \item $\displaystyle a_{n+1}=\frac{pa_n}{qa_n+r}\,,\quad a_{n+1}=\frac{pa_n+q}{ra_n+s}$\\
          \begin{enumerate}
            \item $\displaystyle a_{n+1}=\frac{pa_n}{qa_n+r}$\\
                  逆数型.逆数を取って特性方程式型に帰着.
            \item $\displaystyle a_{n+1}=\frac{pa_n+q}{ra_n+s}$\\
                  一次分数型.$\displaystyle x=\frac{px+q}{rx+s}$を満たす2解$\alpha,\beta$を求めて,\\
                  $\displaystyle b_{n}=\frac{a_{n}-\alpha}{a_{n}-\beta}$と置き,\,$b_{n+1}=\lambda b_n$を満たす$\lambda$を見つけてして等比型に帰着.
          \end{enumerate}
    \item $a_{n+1}+qa_{n+1}+qa_n=0$
          \begin{enumerate}
            \item $a_{n+1}+qa_{n+1}+qa_n=0$\\
                  二次特性方程式型.\,$a_{n+2}+\alpha a_n{n+1}=\beta(a_{n+1}+\alpha a_n{n})$を満たす$\alpha,\beta$を見つけて等比型へ帰着.（$x^2+px+q=0$を解いて$\alpha,\beta$を見つけてもよい.）
            \item $a_{n+1}+qa_{n+1}+qa_n=r$\\
                  等差二次特性方程式型.$b_{n+1}+qb_{n+1}+qa_n=0$となるような$b_n=a_n-\alpha$を満たす$\alpha$を求める.その後は同上.
          \end{enumerate}
    \item 総和型\\
          $S_{n+1}-S_n$をして$a_{n+1}$を生み出す.
    \item 連立型\\
          加減によって得られる$\alpha a_{n}+\beta b_{n}$について,それらが等比数列となるように解を求める.
  \end{enumerate}
  \section{手法の見極め}
    \begin{itemize}
      \item 積のみで表されていたら対数型
      \item 左辺に$a_{n+1}$,右辺に$a_n$を分母に持つ分数が来たら分数型
      \item 困ったら特性方程式を立てて無理やり等比数列に帰着.
      \item 難しそうだと思ったら$a_5$ぐらいまで求めて法則性がないか考える.帰納法も同時に検討.
    \end{itemize}
  \section{式変形のテクニック}
    \begin{itemize}
      \item 左辺を$n+1$,右辺を$n$の式にする.
      \item 漸化式を$n+1$と$n$の関係を表したものから$n+2$と$n+1$を表したものに変えてそれらを足し引きする.
      \item 小さな数で実験する
    \end{itemize}
\end{document}